%%%%%%%%%%%%%%%%%%%%%%%%%%%%%%%%%%%%%%%%%%%%%%%%%%%%%%%%%%%%%%%%%%%%%%%%%%%%%%%
%  Molecular Coding Format manual                by  Akira Yamaji 2026.05.10
%%%%%%%%%%%%%%%%%%%%%%%%%%%%%%%%%%%%%%%%%%%%%%%%%%%%%%%%%%%%%%%%%%%%%%%%%%%%%%%
% ** typeset by LuaLaTeX(luamplib)
% ** mcf2graph must be version 5.30
% ** mcf2graph.mp and main_lib.mcf must be in the same directory to typeset
%----------------------------------------------------------------------------
\documentclass[a4paper]{article}
\usepackage{verbatim}
\usepackage{luamplib}
\usepackage{makeidx}
\usepackage{hyperref}
\makeindex
\hypersetup{colorlinks=true,linkcolor=blue}
%----------------------------------------------------------------------------
\mplibnumbersystem{double}
\mplibcodeinherit{enable}%
\mplibverbatim{enable}%
\mpliblegacybehavior{disabled}%
\begin{mplibcode}
  input mcf2graph;
\end{mplibcode}
%----------------------------------------------------------------------------
\makeatletter%
\edef\fext{mps}%
\topmargin=-18mm
\textheight=254mm
\textwidth=168mm
\oddsidemargin=0mm
%%%%\oddsidemargin=7mm
%%%%\evensidemargin=-7mm
\unitlength=1mm%
%----------------------------------------------------------------------------
\begin{document}
\title{\Huge\sf Molecular Coding Format manual}
\author{Akira Yamaji}
\date{\today}
\maketitle
\begin{center} mcf2graph version 5.30 \end{center}
\begin{center} Located at http://www.ctan.org/pkg/mcf2graph \end{center}
\begin{center} Suggestion or request mail to: mcf2graph@gmail.com \end{center}
%-----------------------------------------------------------------------------
\thispagestyle{empty}
\vspace{5mm}%
\quad
\begin{mplibcode}
max_blength:=10mm;
ratio_thickness_bond:=0.015;
ratio_atom_bond:=0.36;
fmargin:=(2mm,1mm);
fsize:=(40mm,24mm);
blength:=5mm;
sw_frame:=0;
% Glycine
beginfigm drawm("^30,NH2,!2,COOH") endfigm
% L-Alanine
beginfigm drawm("^30,NH2,!~wb,?,!,COOH") endfigm
% L-Valine
beginfigm drawm("^30,NH2,!~wb,/?!,!,COOH") endfigm
% L-Leucine
beginfigm drawm("^30,NH2,!~wb,/!'?!,!,COOH") endfigm
\end{mplibcode}
\\
\begin{mplibcode}
% L-Isoleucine
beginfigm drawm("^30,NH2,!~wb,/?z'!2,!,COOH") endfigm
% L-Serine
beginfigm drawm("^30,NH2,!~wb,/!OH,!,COOH") endfigm
% L-Threonine
beginfigm drawm("^30,NH2,!~wb,/?'!w'OH,!,COOH") endfigm
% L-Cysteine
beginfigm drawm("^30,NH2,!~wb,/!SH,!,COOH") endfigm
\end{mplibcode}
\\
\begin{mplibcode}
% L-Methionine
beginfigm drawm("^30,NH2,!~wb,/!2'S!,!,COOH") endfigm
% L-Phenylalanine
beginfigm drawm("^30,NH2,!~wb,/!Ph,!,COOH") endfigm
% L-Tyrosine
beginfigm drawm("^30,NH2,!~wb,/!Ph'-3:/OH,!,COOH") endfigm
% L-Triptophan
beginfigm drawm("^30,NH2,!~wb,!,COOH,2:,!2,^24,<,?5,-4=Ph,2=dr,5=dl,4:NH")
endfigm
\end{mplibcode}
\\
\begin{mplibcode}
% L-Prorine
beginfigm drawm("^18,?5,3:NH,4:*/COOH") endfigm
% L-Glutamine
beginfigm drawm("^30,NH2,!~wb,!,COOH,2:,!`1,!2,?O,!,NH2") endfigm
% L-Asparagine
beginfigm drawm("^30,NH2,!~wb,/!'?O'!NH2,!,COOH") endfigm
% L-Aspartic acid
beginfigm drawm("^30,NH2,!~wb,/!COOH,!,COOH") endfigm
\end{mplibcode}
\\
\begin{mplibcode}
% L-Glutamic acid
beginfigm drawm("^30,NH2,!~wb,/!2'COOH,!,COOH") endfigm
% L-Lysine
beginfigm drawm("^30,NH2,!~wb,/!4'NH2,!,COOH") endfigm
% L-Arginine
beginfigm drawm("^30,NH2,!~wb,!,COOH,2:,!2`1,!2,NH,!,//NH,!,NH2") endfigm
% L-Hystidine
beginfigm drawm("^30,NH2,!~wb,!,COOH,2:,!2,<,?5,1'3=dl,3:N,5:NH") endfigm
\end{mplibcode}
\\
\begin{mplibcode}
% L-DOPA
beginfigm drawm("^30,NH2,!~wb,!,COOH,2:,!2,Ph,-3'-4:/OH") endfigm
% Ornithine
beginfigm drawm("^30,NH2,!~wb,/!3'NH2,!,COOH") endfigm
% Citrulline
beginfigm drawm("^30,NH2,!~wb,/!3'NH!'?O'!NH2,!,COOH") endfigm
% GABA
beginfigm drawm("^30,NH2,!4,COOH") endfigm
\end{mplibcode}
\\
\begin{mplibcode}
% amino Levulinic acid
beginfigm drawm("^30,NH2,!2,?O,!3,COOH") endfigm
% 4-amino benzoic acid
beginfigm drawm("^30,Ph,1:/NH2,4:/COOH") endfigm
% L-Carboxyl glutamic acid
beginfigm drawm("^30,NH2,!~wb,/!'?COOH'!COOH,!,COOH") endfigm
% L-Hydroxy Prorine
beginfigm drawm("^18,?5,1:/OH,3:NH,4:*/COOH") endfigm
\end{mplibcode}
%-----------------------------------------------------------------------------
\twocolumn
\thispagestyle{empty}
\tableofcontents
\linethickness{0.08mm}%
%----------------------------------------------------------------------------
\newpage
\twocolumn
\section{Introduction}
Molecular Coding Format(MCF) is new linear notation
represent chemical structure diagrams.
This Coding is named from programing technique
such as operator, array, scope, macro, adressing, etc.
mcf2graph convert from MCF to PNG, SVG, MOL file.
It is also able to calculate molecular weight,
exact mass, molecular formula.
%-----------------------------------------------------------------------------
\section{MCF syntax}
\subsection{Make bond}
\subsubsection{Chain}
\begin{verbatim}
real number plus (+): counterclockwize 
real number minus(-): clockwize
$n (0<=n<360): absolute angle

^10,-30,45,-45,60,$300,$0
\end{verbatim}
%-----------------------------------------
\begin{mplibcode}
sw_frame:=Outside;
blength:=0;
beginfigm
  #(60mm,18mm);
  sw_numbering:=Bond;
  ratio_chain_ring:=1;
%----------------------------------------------------------------------
  drawm("^15,-30,45,-45,60,$300,$0,1^$15'2^$345'3^$30'4^$345:?~dt,5'6=vf")
  ++(
    defaultscale:=0.5;
    labeloffset:=2bp;
    drawarrow B7/*.7{B7right}..{B1right}B1/*.7;   label.ulft("-30",B7/*.7);
    drawarrow B8/*.7{B8left}..{B2left}B2/*.7;     label.llft("45",B8/*.7);
    drawarrow B9/*.7{B9right}..{B3right}B3/*.7;   label.ulft("-45",B9/*.7);
    drawarrow B10/*.7{B10left}..{B4left}B4/*.7;   label.llft("60",B10/*.7);
  )
endfigm
\end{mplibcode}
%-----------------------------------------------------------------------------
\subsubsection{Chain with !,!n}
\index{"!}%
\begin{verbatim}
!  : take value 60 or -60 depend on
     current angle
!6 : !,!,!,!,!,!

^-30,!6
\end{verbatim}
%-----------------------------------------
\begin{mplibcode}
beginfigm
  #(60mm,20mm);
  sw_numbering:=Bond;
  ratio_chain_ring:=1;
%----------------------------------------------------------------------
  drawm("^-30,!6,1^-120'2^60'3^-60'4^60'5^-60'6^60:?~dt")
  ++(
    defaultscale:=0.5;
    labeloffset:=2bp;
    drawarrow B7/*.7{B7right}..{B1right}B1/*.7;   label.ulft("-60",B7/*.7);
    drawarrow B8/*.7{B8left}..{B2left}B2/*.7;     label.llft("60",B8/*.7);
    drawarrow B9/*.7{B7right}..{B3right}B3/*.7;   label.ulft("-60",B9/*.7);
    drawarrow B10/*.7{B10left}..{B4left}B4/*.7;   label.llft("60",B10/*.7);
    drawarrow B11/*.7{B11right}..{B5right}B5/*.7; label.ulft("-60",B11/*.7);
    drawarrow B12/*.7{B12left}..{B6left}B6/*.7;   label.llft("60",B12/*.7);
  )
endfigm
\end{mplibcode}
%-----------------------------------------------------------------------------
\begin{verbatim}
take value 0 just after jump to atom
n:,! : !=0

^-30,!5,3:,!
\end{verbatim}
%-----------------------------------------
\begin{mplibcode}
beginfigm
  #(60mm,16mm) #@(2mm,2mm)
  sw_trimming:=1;
  sw_numbering:=Atom;
  ratio_chain_ring:=1;
  drawm("^-30,!6,3:,!")
endfigm
\end{mplibcode}
%-----------------------------------------------------------------------------
\subsubsection{Jump to atom}
\begin{verbatim}
n: : Jump to An
** An: atom number(-999<=n<=4095)

^-30,!6,3:,0,5:,-30
\end{verbatim}
%-----------------------------------------
\begin{mplibcode}
beginfigm
  #(60mm,16mm) #@(2mm,2mm)
  sw_trimming:=1;
  sw_numbering:=Atom;
  ratio_chain_ring:=1;
  drawm("^-30,!6,3:,0,5:,-30")
endfigm
\end{mplibcode}
%-----------------------------------------------------------------------------
\subsubsection{Branch bond}
\index{\textasciicircum}% ^
\index{\textasciitilde}% ~
\index{`}%
\index{"!w}%
\index{"!z}%
\index{"!d}%
\begin{verbatim}
n:,!   : n:,!
n:,!w  : n:,!~wf
n:,!z  : n:,!~zf
n:,!d  : n:,!~db
n:,!dr : n:,!~dr
n:,!t  : n:,!~tm

^30,!12,
2:,!,4:,!w,6:,!z,8:,!d,10:,!dr,12:,!t
\end{verbatim}
%-----------------------------------------
\begin{mplibcode}
beginfigm
  #(80mm,20mm) #@(2mm,2mm)
  sw_numbering:=Atom;
  ratio_chain_ring:=1;
  drawm("^30,!12,2:,!,4:,!w,6:,!z,8:,!d,10:,!dr,12:,!t")
endfigm
\end{mplibcode}
%-----------------------------------------------------------------------------
\subsubsection{Rotate current angle}
\index{\textasciicircum}%
\begin{verbatim}
^angle : rotate current angle

0,0,^90,0,^-90,0,^$315,0,^$90,0,^$0,0 
\end{verbatim}
%-----------------------------------------
\begin{mplibcode}
beginfigm
  #(62mm,22mm) #@(2mm,3mm)
  sw_trimming:=1;
  sw_numbering:=Bond;
  ratio_chain_ring:=1;
    defaultscale:=0.5;
    labeloffset:=2bp;
  drawm("`1,0,0,^90,0,^-90,0,^$315,0,^$90,0,^$0,0,1_7=vf,3'4^180:?~dt")
  ++(
    drawarrow B8/*.7{B8left}..{B3left}B3/*.7;   label.urt("90",B8/*.7);
    drawarrow B9/*.7{B9right}..{B4right}B4/*.7; label.urt("-90",B9/*.7);
  )
endfigm
\end{mplibcode}
%-----------------------------------------------------------------------------
\subsubsection{Connect atom}
\begin{verbatim}
:n : Connect to An

^-30,!6,3:,!,!3,:6~bd,9:,:4~bz
\end{verbatim}
%-----------------------------------------
\begin{mplibcode}
beginfigm
  #(60mm,20mm) #@(2mm,2mm)
  sw_trimming:=1;
  sw_numbering:=Atom;
  ratio_chain_ring:=1;
  %----------------------------------
  drawm("^-30,!6,3:,!,!3,:6~bd,9:,:4~bz")
endfigm
\end{mplibcode}
%-----------------------------------------------------------------------------
\subsubsection{Ring}
\index{?n}%
\begin{verbatim}
?n : n membered ring(3<=n<=20)
?6
\end{verbatim}
%-----------------------------------------
\begin{mplibcode}
beginfigm
  #(60mm,20mm) #@(2mm,3mm)
  sw_trimming:=1;
  sw_numbering:=Bond;
  %---------------------------
  drawm("?6")
endfigm
\end{mplibcode}
%-----------------------------------------------------------------------------
\subsubsection{Ring length}
\begin{verbatim}
?n`length : change ring length

?6,4:,!,?6`1.2
\end{verbatim}
%-----------------------------------------
\begin{mplibcode}
beginfigm
  #(60mm,20mm) #@(2mm,2mm)
  sw_trimming:=1;
  sw_numbering:=Bond;
  drawm("?6,4:,!,?6`1.2")
endfigm
\end{mplibcode}
%-----------------------------------------------------------------------------
\subsection{Change bond type}
\subsubsection{Double,triple,wedge,vector}
\index{\textasciitilde}%
\index{\textasciitilde\textasciitilde}%
\index{dm}%
\index{dl}%
\index{dr}%
\index{db}%
\index{tm}%
\index{"!db}%
\index{"!tm}%
\index{"!d}%
\index{"!t}%
\index{"!dm}%
\index{"!dl}%
\index{"!dr}%
\begin{verbatim}
(Double,triple)
dm  : double middle
dl  : double left side
dr  : double right side
db  : double left or right side
tm  : triple
!db,!d  : !~db  / !tm,!t : !~tm

^-30,!~dm,!,!~dl,!,!~dr,!~db,!~db,!,!~tm
^-30,!dm ,!,!dl ,!,!dr ,!d  ,!d  ,!,!t
\end{verbatim}
%-----------------------------------------
\begin{mplibcode}
beginfigm
  #(70mm,10mm) #@(2mm,2mm)
  sw_trimming:=0;
  ratio_chain_ring:=1;
  drawm("^-30,!~dm,!,!~dl,!,!~dr,!~db,!~db,!,!~tm")
  ++(defaultscale:=0.6; labeloffset:=0;
      label.rt("(dm)",A1+(0,-0.7l));
      label.rt("(dl)",A3+(0,-0.7l));
      label.rt("(dr)",A5+(0,-0.7l));
      label.rt("(db)",A6+(0,-0.2l));
      label.rt("(db)",A7+(0,-0.7l));
      label.rt("(tm)",A9+(0,-0.7l));
  )
endfigm
\end{mplibcode}
\vspace{-3mm}%
%-----------------------------------------------------------------------------
%%\subsubsection{Wedge}
\index{wf}%
\index{wb}%
\index{zf}%
\index{zb}%
\index{vf}%
\index{vb}%
\begin{verbatim}

(Wedge,Vector)
wf: wedge forward
wb: wedge backward
zf: hashed(zebra stripe) wedge foward
zb: hashed(zebra stripe) wedge backward
vf:vector forward  /  vb:vector backward

^-30,
 !~wf,!,!~wb,!,!~zf,!,!~zb,!,!~vf,!~vb
\end{verbatim}
%-----------------------------------------
\begin{mplibcode}
beginfigm
  #(70mm,10mm) #@(2mm,2mm)
  sw_trimming:=0;
  ratio_chain_ring:=1;
  drawm("^-30,!~wf,!,!~wb,!,!~zf,!,!~zb,!,!~vf,!,!~vb")
  ++(defaultscale:=0.6; labeloffset:=0;
      label.rt("(wf)",A1+(0,-0.7l));
      label.rt("(wb)",A3+(0,-0.7l));
      label.rt("(zf)",A5+(0,-0.7l));
      label.rt("(zb)",A7+(0,-0.7l));
      label.rt("(vf)",A9+(0,-0.7l));
      label.rt("(vb)",A11+(0,-0.7l));
  )
endfigm
\end{mplibcode}
\vspace{-3mm}%
%-----------------------------------------------------------------------------
%%\subsubsection{Dotted,wave}
\index{dt}%
\index{wv}%
\index{bd}%
\index{bz}%
\begin{verbatim}

(Dotted,wave)
Bn=bond type : change bond type at Bn
dt : dotted   /  wv : wave
bd : broad    /  bz : broad dotted 

^-30,!7,1=dt,3=wv,5=bd,7=bz
\end{verbatim}
%-----------------------------------------
\begin{mplibcode}
beginfigm
  #(70mm,10mm) #@(2mm,2mm) |=(9mm)
  sw_trimming:=1;
  ratio_chain_ring:=1;
  drawm("^-30,!7,1=dt,3=wv,5=bd,7=bz")
  ++(defaultscale:=0.6; labeloffset:=0;
      label.rt("(dt)",A1+(0,-0.6l));
      label.rt("(wv)",A3+(0,-0.6l));
      label.rt("(bd)",A5+(0,-0.6l));
      label.rt("(bz)",A7+(0,-0.6l));
  )
endfigm
\end{mplibcode}
\vspace{-3mm}%
%-----------------------------------------------------------------------------
\subsubsection{Over line}
\index{si\_}%
\index{wf\_}%
\index{wb\_}%
\index{zf\_}%
\index{zb\_}%
\index{bd\_}%
\index{dl\_}%
\index{dr\_}%
\index{dm\_}%
\begin{verbatim}
si_ : single over line 
wf_ : wedge forward over line 
wb_ : wedge backward over line 
zf_ : hashed wedge forward over line 
zb_ : hashed wedge backward over line 
bd_ : broad over line 
dl_ : duble left over line 
dr_ : duble right over line 
dm_ : duble over line 

^30,!8,!,60,90`18,
 2~si_'4~wf_'6~wb_'8~zf_'10~zb_'
  12~bd_'14~dl_'16~dr_'18~dm_:?`2
\end{verbatim}
%-----------------------------------------
\begin{mplibcode}
beginfigm
  sw_trimming:=1;
  #(75mm,20mm);
  ratio_chain_ring:=1;
  drawm("^30,!18,$90`1.5,90`15.5,"&
         "2~si_'4~wf_'6~wb_'8~zf_'10~zb_'"&
         "12~bd_'14~dl_'16~dr_'18~dm_:?`2")
endfigm
\end{mplibcode}
%-----------------------------------------------------------------------------
\subsubsection{Steric ring}
\index{wf\_r}%
\index{wb\_r}%
\index{bd\_r}%
\begin{verbatim}
wf_r : wedge foward (half width)
bd_r : broad (half width, rounded)
wb_r : wedge backward (half width)

`1.25,-30~wf_r,30~bd_r`1,30~wb_r,
 120,O,30,:1,`,`.5,6^$90:/!OH,
 1^$270'2^$90'3^$270'4^$90:/OH
\end{verbatim}
%-----------------------------------------
\begin{mplibcode}
beginfigm
  #(75mm,16mm);
  @(0.05,0.5)
  drawm("`1.25,-30~wf_r,30~bd_r`1,30~wb_r,120,O,30,:1,`,"&
        "`.5,1^$270'2^$90'3^$270'4^$90:/OH,6^$90:/!OH")
  defaultscale:=0.6;
  @(0.5,0.7)  drawm("0~wf_r") ++(label.lft("wf_r:",A1);)
  @(0.5,0.2)  drawm("0~wf")   ++(label.lft("wf:",A1);)
  @(0.75,0.7) drawm("0~bd_r") ++(label.lft("bd_r:",A1);)
  @(0.75,0.2) drawm("0~bd")   ++(label.lft("bd:",A1);)
  @(1,0.7)    drawm("0~wb_r") ++(label.lft("wb_r:",A1);)
  @(1,0.2)    drawm("0~wb")   ++(label.lft("wb:",A1);)
endfigm
\end{mplibcode}
%-----------------------------------------------------------------------------
\subsubsection{Change multiple bond type}
\index{vf}%
\begin{verbatim}
2'4'6=dl or 2'4'6=dl : 2=dl,4=dl,6=dl

^30,!7,2'4'6=dl or ^30,!7,2'4'6=dl 
\end{verbatim}
%-----------------------------------------
\begin{mplibcode}
beginfigm
  #(60mm,8mm) #@(2mm,3mm)
  ratio_chain_ring:=1;
  drawm("^30,!7,2'4'6=dl")
endfigm
\end{mplibcode}
%-----------------------------------------------------------------------------
\subsection{Change bond length}
\subsubsection{Chanege chain length}
\index{`}%
\begin{verbatim}
(!,!n)`length : change length of !,!n

^-30,!2,!4`1.2,!2
\end{verbatim}
%-----------------------------------------
\begin{mplibcode}
beginfigm
  #(55mm,8mm);
  sw_numbering:=Bond;
  drawm("^-30,!2,!4`1.2,!2")
endfigm
\end{mplibcode}
%-----------------------------------------------------------------------------
\subsubsection{Change multiple bond length}
\begin{verbatim}
`n : bond length=n
` : reset bond length

^-30,!2,`1.2,!4,`,!2
\end{verbatim}
%-----------------------------------------
\begin{mplibcode}
beginfigm
  #(55mm,8mm);
  sw_numbering:=Bond;
  drawm("^-30,!2,`1.2,!4,`,!2")
endfigm
\end{mplibcode}
%-----------------------------------------------------------------------------
\subsection{Change atom}
\subsubsection{Insert atom}
\begin{verbatim}
Insert hetero atom

^-30,!2,O,!2,N,!2
\end{verbatim}
%-----------------------------------------
\begin{mplibcode}
beginfigm
  sw_trimming:=1;
  #(50mm,7mm);
  drawm("^-30,!2,O,!2,N,!2")
endfigm
\end{mplibcode}
%-----------------------------------------------------------------------------
\subsubsection{Change atom}
\index{:}%
\begin{verbatim}
2:O : change A2 C to O
3'4:N : change A3,A4 C to N

^30,!4,2:O,3'4:N
\end{verbatim}
%-----------------------------------------
\begin{mplibcode}
beginfigm
  #(70mm,10mm) ##(0.48,1)
  @(0,0.5) drawm("^30,!5,2:O,3'4:N")
  sw_numbering:=Atom;
  @(1,0.5) drawm("^30,!5,2:O,3'4:N")
endfigm
\end{mplibcode}
%-----------------------------------------------------------------------------
\subsubsection{Brock address}
\begin{verbatim}
< : divide brock

?6,4:,!,<,?6,2:O
\end{verbatim}
%-----------------------------------------
\begin{mplibcode}
beginfigm
  #(70mm,14mm) #@(3mm,1.5mm)
  @(0,.5) drawm("?6,4:,!,<,?6,2:O")
  sw_numbering:=Atom;
  ##(1,.88)
  @(1,.5) drawm("?6,4:,!,<,?6,2:O")
endfigm
\end{mplibcode}
%-----------------------------------------------------------------------------
\subsubsection{Reset brock address}
\begin{verbatim}
> : reset brock adress

?6,4:,!,<,?6,>,2:N
\end{verbatim}
%-----------------------------------------
\begin{mplibcode}
beginfigm
  #(70mm,14mm) #@(3mm,1.5mm)
  @(0,.5) drawm("?6,4:,!,<,?6,>,2:N")
  sw_numbering:=Atom;
  ##(1,.88)
  @(1,.5) drawm("?6,4:,!,<,?6,>,2:N")
endfigm
\end{mplibcode}
%-----------------------------------------------------------------------------
\subsubsection{Absolute address}
\index{\$}% $
\begin{verbatim}
$2:n : change A$2 C to N
**1<=n<=3095

?6,4:,!,<,?6,$2:N
\end{verbatim}
%-----------------------------------------
\begin{mplibcode}
beginfigm
  #(70mm,14mm) #@(3mm,1.5mm)
  @(0,.5) drawm("?6,4:,!,?6,$2:N")
  sw_numbering:=Atom;
  ##(1,.88) 
  @(1,.5) drawm("?6,4:,!,?6,$2:N")
endfigm
\end{mplibcode}
%-----------------------------------------------------------------------------
\subsubsection{Relative address}
\begin{verbatim}
-2:n : change A(-2) C to N
**-999<=n<=-1

?6,4:,!,?6,-2:N
\end{verbatim}
%-----------------------------------------
\begin{mplibcode}
beginfigm
  #(70mm,14mm) #@(3mm,1.5mm)
  @(0,.5) drawm("?6,4:,!,?6,-2:N")
  sw_numbering:=Atom;
  ##(1,.88)
  @(1,.5) drawm("?6,4:,!,?6,-2:N")
endfigm
\end{mplibcode}
%-----------------------------------------------------------------------------
\subsubsection{Charged atom}
\begin{verbatim}
p : positive
n : negative

^-30,!2,N,??,p,!2,S,n^180,
!6,7:N,7:??,9:S,7:n,9:n^180
\end{verbatim}
%-----------------------------------------
\begin{mplibcode}
beginfigm
  sw_trimming:=1;
  #(60mm,12mm);
  drawm("^-30,!2,N,??,p,!2,S,n^180,!6,7:N,7:??,9:S,7:p,9:n^180")
endfigm
\end{mplibcode}
%-----------------------------------------------------------------------------
\newpage
%-----------------------------------------------------------------------------
\subsection{Fuse ring}
%%%\subsubsection{Attached 1 bond}
\begin{verbatim}
(Attached 1 bond)

?6,3=?6 : fuse ?6 at B3
** Bn(n:-999<=n<=4095): bond number

?6,3=?6
\end{verbatim}
%-----------------------------------------
\begin{mplibcode}
beginfigm
  #(60mm,18mm) #@(2mm,1.5mm)
  sw_trimming:=1;
  sw_numbering:=Bond;
  drawm("^30,?6,3=?6,3=dt,7_11=bd_r")
endfigm
\end{mplibcode}
%-----------------------------------------------------------------------------
\begin{verbatim}
** fused ring size depend on 
attached bond length

?6,4:,!,?6`1.2,5=?6,11=?6
\end{verbatim}
%-----------------------------------------
\begin{mplibcode}
beginfigm
  #(60mm,25mm) #@(2mm,2mm)
  sw_trimming:=1;
  sw_numbering:=Bond;
  drawm("?6,4:,!,?6`1.2,5=?6,11=?6,14_23=bd_r,5'11=dt")
endfigm
\end{mplibcode}
%-----------------------------------------------------------------------------
\begin{verbatim}
?6,3=`1.3'?6 : fuse #1.3'?6 at B3
`1.3'?6 : 6 membered ring scaled 1.3
** `n'?m (5<=m<=6,1.1<=n<=1.5)

?6,3=`1.3'?6
\end{verbatim}
%-----------------------------------------
\begin{mplibcode}
beginfigm
  #(60mm,18mm) #@(2mm,1.5mm)
  margin_top_bottom:=1.5mm;
  sw_numbering:=Bond;
  sw_trimming:=1;
  drawm("^30,?6,3=`1.3'?6,3=dt,7_11=bd_r")
endfigm
\end{mplibcode}
%-----------------------------------------------------------------------------
\begin{verbatim}
?6,-3'-4'-4'-2'-2'-4'-4=?6
?6,4'8'13'20'25'28'33=?6
\end{verbatim}
%-----------------------------------------
\begin{mplibcode}
beginfigm
  #(70mm,20mm) #@(2mm,2mm)
  sw_numbering:=Bond;
  sw_trimming:=1;
  drawm("^30,?6,-3'-4'-4'-2'-2'-4'-4=?6,4'8'13'20'25'28'33=dt")
endfigm
\end{mplibcode}
%-----------------------------------------------------------------------------
\index{--}%
\begin{verbatim}
(Attached 2 bond)

4--11=?6 : fuse 4/6 ring to B11..B4
4--11=?5 : fuse 3/5 ring to B11..B4
4--11=?4 : fuse 2/4 ring to B11..B4

1:^30,?6,3=?6,11--4=?6
2:^30,?6,3=?6,11--4=?5
3:^30,?6,3=?6,11--4=?4
\end{verbatim}
%-----------------------------------------
\begin{mplibcode}
beginfigm
  #(75mm,20mm) #@(2mm,2mm)
  sw_numbering:=Bond;
  ##(1,.9)
  @( 0,.5) drawm("^30,?6,3'11--4=?6,11'4=dt,12_15=bd_r")
  ++(defaultscale:=0.4; label("(1)",p0);)
  ##(1,.9)
  @(.5,.5) drawm("^30,?6,3=?6,11--4=?5,11'4=dt,12_14=bd_r")
  ++(defaultscale:=0.4; label("(2)",p0);)
  ##(1,.9)
  @( 1,.5) drawm("^30,?6,3=?6,11--4=?4,11'4=dt,12'13=bd_r")
  ++(defaultscale:=0.4; label("(3)",p0);)
endfigm
\end{mplibcode}
%%%%%%%\vspace{-3mm}%
%-----------------------------------------------------------------------------
\index{---}%
\begin{verbatim}
(Attached 3 bond)

16---4=?6 : fuse 3/6 ring to B16..B4
16---4=?5 : fuse 2/5 ring to B16..B4

1:?6,3'10'16---4=?6
2:?6,3'10=?6,16---4=?5

\end{verbatim}
%-----------------------------------------
\begin{mplibcode}
beginfigm
  #(60mm,20mm) #@(2mm,2mm)
  sw_numbering:=Bond;
  @(0,1) drawm("?6,3'10=?6,16---4=?6,16'4=dt,17_19=bd_r")
  ++(defaultscale:=0.4; label("(1)",p0);)
  @(1,0) drawm("?6,3'10=?6,16---4=?5,16'4=dt,17'18=bd_r")
  ++(defaultscale:=0.4; label("(2)",p0);)
endfigm
\end{mplibcode}
%%%%%%%%\vspace{-3mm}%
%-----------------------------------------------------------------------------
\index{----}%
\begin{verbatim}
(Attached 4 bond)

21----4=?6 : fuse 2/6 ring to B21..B4

^-30,?6,3'10'15'(21----4)=?6

\end{verbatim}
%-----------------------------------------
\begin{mplibcode}
beginfigm
  #(60mm,20mm) #@(2mm,2mm)
  sw_numbering:=Bond;
  drawm("^-30,?6,3'10'15'21----4=?6,21'4=dt,22'23=bd_r")
endfigm
\end{mplibcode}
%-----------------------------------------------------------------------------
\subsection{Spiro ring}
\begin{verbatim}
:4,?5 : add ?5 at A4

^30,!6,4:,?5
\end{verbatim}
%-----------------------------------------
\begin{mplibcode}
beginfigm
  #(40mm,15mm);
  sw_numbering:=Atom;
  ratio_chain_ring:=1;
  drawm("^30,!6,4:,?5")
endfigm
\end{mplibcode}
%-----------------------------------------------------------------------------
\subsection{Group}
\subsubsection{Insert group}
\index{?}%
\index{methyl}%
\index{??}%
\index{dimethyl}%
\index{?O}%
\index{carbonyl}%
\index{?S}%
\index{thioketone}%
\index{?NH}%
\index{imino}%
\index{ethyl}%
\index{propyl}%
\index{isopropyl}%
\index{tert-butyl}%
\index{/}%
\index{Ph}%
\index{phenyl}%
\begin{verbatim}
?    : methyl
??   : dimethyl
?O   : carbonyl
?S   : thioketone
?NH  : imino
/!   : ethyl
/!2  : propyl
/?!  : isopropyl
/??! : tert-butyl
/Ph  : phenyl

^30,!,?,!2,??,!2,?O,!2,?S,!2,
 ?NH,!2,/!,!2,/!2,!2`1,/?!,!2`1,
 /??!,!2`1,/`.6'Ph,!
\end{verbatim}
%-----------------------------------------
\begin{mplibcode}
beginfigm
  #(70mm,20mm);
  drawm("^30,!,?,!2,??,!2,?O,!2,?S,!2,?NH,!2,/!,!2,/!2,!2`1,"&
        "/?!,!2`1,/??!,!2`1,/`.6'Ph,!")
endfigm
\end{mplibcode}
%-----------------------------------------------------------------------------
\subsubsection{Insert modified group}
\index{//}%
\index{double}%
\index{*/}%
\index{wedge forward}%
\index{/*}%
\index{hashed wedge}%
\index{*/*}%
\index{wave}%
\index{direct}%
\begin{verbatim}
//  : double (double middle)
*/  : wedge forward
/*  : hashed wedge forward
*/* : wave

^30,!,//!,!2,*/H,!2,/*H,!2,*/*H,!
\end{verbatim}
%-----------------------------------------
\begin{mplibcode}
beginfigm
  #(70mm,14mm);
  drawm("^30,!`1,//!,!2`1,*/H,!2`1,/*H,!2`1,*/*H,!`1")
  ++(defaultscale:=0.75;
      label("//",A2-(0,0.45l));
      label("*/",A4-(0,0.45l));
      label("/*",A6-(0,0.45l));
      label("*/*",A8-(0,0.45l));
     )
endfigm
\end{mplibcode}
%-----------------------------------------------------------------------------
\index{\textasciicircum}% ^
\index{\textasciitilde}% ~
\index{`}%
\index{\textless}%
\index{change type}%
\index{change angle}%
\index{change length}%
\index{change environment}%
\begin{verbatim}
~ : change type
^ : change angle,environment
` : change length

^-30,`1,!,
  ?`2^30,!2,/!2^lr,!2,/!2^rl,!
\end{verbatim}
%-----------------------------------------
\begin{mplibcode}
beginfigm
  #(60mm,16mm);
  sw_trimming:=1;
  drawm("^-30,`1,!2,?`2^30,!2,/!2^lr,!2,/!2^rl,!")
endfigm
\end{mplibcode}
%-----------------------------------------------------------------------------
\subsubsection{Add group}
\index{?w}%
\index{?z}%
\index{?d}%
\begin{verbatim}
?w : ?~wf
?z : ?~zf
?d : ?~dm

^30,!17,2:?,3:/!,4:/!2,7:/?!,
 8:/??!,10^-15:/`.6'Ph,
  11'12':?w,13'14:?z,15'16:?d
\end{verbatim}
%-----------------------------------------
\begin{mplibcode}
beginfigm
  #(75mm,25mm) #@(2mm,2mm)
  sw_numbering:=Atom;
  drawm("^30,!17,2:?,3:/!,4:/!2,7:/?!,8:/??!,10^-15:/`.6'Ph,"&
     "11'12:?w,13'14:?z,15'16:?d")
endfigm
\end{mplibcode}
%-----------------------------------------------------------------------------
\subsubsection{Add modified group}
\begin{verbatim}
~,^,` : change type,angle,length

^30,!6,2~wf'4~zf'6^-30'8^$120:?
\end{verbatim}
%-----------------------------------------
\begin{mplibcode}
beginfigm
  #(60mm,18mm) #@(2mm,2mm)
  sw_trimming:=1;
  sw_numbering:=Atom;
  drawm("^30,!8`1,2~wf'4~zf'6^-30'8^$120:?")
endfigm
\end{mplibcode}
%-----------------------------------------------------------------------------
\begin{verbatim}

^,` : change angle,length,environment

^-30,!7`1,3:?`2^30,5:/!2^lr,7:/!2^rl
\end{verbatim}
%-----------------------------------------
\begin{mplibcode}
beginfigm
  #(60mm,16mm);
  sw_trimming:=1;
  sw_numbering:=Atom;
  drawm("^-30,!7`1,3:?`2^30,5:/!2^lr,7:/!2^rl")
endfigm
\end{mplibcode}
%-----------------------------------------------------------------------------
\newpage
%-----------------------------------------------------------------------------
\subsection{Chain environment}
\subsubsection{Horizontal,vertical}
\index{hz}%
\index{vt}%
\index{"'}%
\begin{verbatim}
^hz : horizontal environment (default)
^vt : vertical environment

?4,
 3^-90'3^-30'3^90:/!3^hz,
 1^-60'1'1^60:/!3^vt
\end{verbatim}
%-----------------------------------------
\begin{mplibcode}
beginfigm
  #(60mm,35mm);
  sw_trimming:=1;
  ratio_chain_ring:=1;
  drawm("?4,3^-90'3^-30'3^90:/!3^hz,1^-60'1'1^60:/!3^vt")
  ++(defaultscale:=0.5; labeloffset:=2bp;
      label.rt("^hz",A8);
      label.top("^hz",A12);
      label.top("^hz",A16);
      label.rt("^vt",A20);
      label.top("^vt",A24);
      label.rt("^vt",A28);
  )
endfigm
\end{mplibcode}
%-----------------------------------------------------------------------------
\subsubsection{Left-right,right-left}
\index{lr}%
\index{rl}%
\begin{verbatim}
^lr : left-right environment
^rl : right-left environment

^-30,!6,
3^-30'3'3^30:/!3^lr,
\end{verbatim}
%-----------------------------------------
\begin{mplibcode}
beginfigm
%%%%%  EN:Left-right_right-left
  #(50mm,24mm);
%  sw_trimming:=1;
  drawm("^-30,!6,3^-30'3'3^30:/!3^lr,5^-30'5'5^30:/!3^rl")
  ++(defaultscale:=0.5; labeloffset:=2bp;
      label.lft("^lr",A11);
      label.lft("^lr",A15);
      label.lft("^lr",A19);
      label.rt("^rl",A23);
      label.rt("^rl",A27);
      label.rt("^rl",A31);
  )
endfigm
\end{mplibcode}
%-----------------------------------------------------------------------------
\subsubsection{Rotate fixed angle}
\index{\textasciicircum}%
\begin{verbatim}
^*n : rotate n

^30,!4,
2:/!6^*30,  % :2,!,30,30,30,30,30,30
4:/!4^*-45  % :4,!,-45,-45,-45,-45

\end{verbatim}
%-----------------------------------------
\begin{mplibcode}
beginfigm
  #(45mm,24mm);
  sw_trimming:=1;
  drawm("^30,!4,2:/!6^*30,4:/!4^*-45")
  ++(defaultscale:=0.5; labeloffset:=2bp;
      label.bot("^*30",A12);
      label.bot("^*-45",A17);
  )
endfigm
\end{mplibcode}
%-----------------------------------------------------------------------------
\subsubsection{Rotate multiple angle}
\begin{verbatim}
^(90'-90,...) : rotate 90,-90,...

^30,!6,6:/!5^(90'-90'90'-90'90)
\end{verbatim}
%-----------------------------------------
\begin{mplibcode}
beginfigm
  #(62mm,24mm);
  drawm("^30,!6,6:/!5^(90'-90'90'-90'90)")
  ++(defaultscale:=0.5; labeloffset:=2bp;
      label.lft("^(90'-90'90'-90'90)",A13);
      label.rt("^*90",A12);
      label.lft("^*-90",A11);
      label.rt("^*90",A10);
      label.lft("^*-90",A9);
      label.rt("^*90",A8);
  )
endfigm
\end{mplibcode}
%-----------------------------------------------------------------------------
\newpage
\subsection{Miscellaneous}
%-----------------------------------------------------------------------------
\subsubsection{Abbreviated parts}
\index{NH}%
\index{N"!}%
\index{N"!2}%
\index{S?O}%
\index{S?O?O}%
\begin{verbatim}
NH : N,/H~nl   N? : N,? 
N?2 : N,/!     S?O : S,?O
S?O?O : S,?O^35,//^-35

^-30,!2,NH,!2,N?,!2,N?2,S?O,!2,S?O?O,!
\end{verbatim}
%-----------------------------------------
\begin{mplibcode}
beginfigm
#(60mm,10mm);
sw_trimming:=1;
drawm("^-30,!2`1.2,NH,!2`1.2,N?,!2`1.2,"&
       "N?2,!2`1.2,S?O,!2`1.2,S?O?O,!`1.2")
endfigm
\end{mplibcode}
%-----------------------------------------------------------------------------
\index{?"!}%
\index{??}%
\index{??"!}%
\index{N?"!}%
\index{dimethylamino}%
\begin{verbatim}
??   : ?^35,?-35  /?!  : isopropyl
/??! : tert-butyl   /N?! : dimethylamino

^30,!11`1,2:??,4:/!,6:/?!,8:/??!,10:/N?!
\end{verbatim}
%-----------------------------------------
\begin{mplibcode}
beginfigm
#(60mm,12mm);
sw_trimming:=1;
drawm("^30,!11`1.2,2:??,4:/!,6:/?!,8:/??!,10:/N?!")
endfigm
\end{mplibcode}
%-----------------------------------------------------------------------------
\subsubsection{Define group,parts}
\index{define group}%
\index{define group with atoms}%
\index{define parts}%
\index{`}%
\index{``}%
\begin{verbatim}
`(..) : define group
``(..)(..) : define group with atoms
'(..) : define parts
''("..") : define parts with string

CF3:=`("{CF3}"); % ** group weight =0
CF3:=``("{CF3}")(C,F,F,F);
  ** group weight =12(C)+19(F)*3=69

iBuOH:='(!?!,OH);
drawm("^30,?6,4'6:/iBuOH,3:/CF3")
\end{verbatim}
%-----------------------------------------
\begin{mplibcode}
beginfigm
  #(60mm,14mm);
  sw_trimming:=1;
  iBuOH:='(!?!,OH);
  CF3:=``("{CF3}")(C,F,F,F);
  drawm("^30,?6,4'6:/iBuOH,3:/CF3")
endfigm
\end{mplibcode}
%-----------------------------------------------------------------------------
\subsubsection{Concatenate group,parts}
\index{'}%
\begin{verbatim}
a'b : concat a,b

^30,NH2,!wb,/!Ph'5:/OH,!COOH
\end{verbatim}
%-----------------------------------------
\begin{mplibcode}
beginfigm
  #(60mm,14mm) #@(2mm,1mm)
  sw_trimming:=1;
  drawm("^30,NH2,!wb,/!Ph'5:/OH,!COOH")
endfigm
\end{mplibcode}
%-----------------------------------------------------------------------------
\subsubsection{Move position [ (): ]}
\begin{verbatim}
 (x'y): : Move l*(x,y) from current position
$(x'y): : Move l*(x,y) from origin(@1)
          ** l=bond length of ring

^30,?6,3:,!4,?O,!,O,n_^60,$(6'1):,H,p^15
\end{verbatim}
%-----------------------------------------
\begin{mplibcode}
beginfigm
  #(70mm,15mm) #@(2mm,1mm)
  sw_trimming:=1;
  drawm("^30,?6,3:,!,!3,?O,!,O,n^60,$(6'1):,H,p^15")
  ++(drawdot A1 withpen pencircle scaled 3bp;
     drawdot A12 withpen pencircle scaled 3bp;
      pickup pencircle scaled 0.1bp;
      for i=0 upto 6:
        draw (A1+(l*i,l-4bp))--(l*i,l+4bp);
      endfor
      draw A1--(A1+(0,1l))--A1+(6l,1l);
      draw (A1+(-4bp,1l))--(A1+(4bp,1l));
      draw (A1+(-4bp,0))--(A1+(4bp,0));
  )
endfigm
\end{mplibcode}
%-----------------------------------------------------------------------------
\subsubsection{Serial number}
\begin{verbatim}
6_10 : 6,7,8,9,10
^30,!14,2'6_10'14:?~bd_r`0.5
\end{verbatim}
%-----------------------------------------
\begin{mplibcode}
beginfigm
  #(75mm,12mm) |<(8mm)
  sw_numbering:=Atom;
  drawm("^30,!14,2'6_10'14:?~bd_r`0.5")
endfigm
\end{mplibcode}
%-----------------------------------------------------------------------------
\subsubsection{Change color}
\index{red}%
\index{blue}%
\index{green}%
\begin{verbatim}
beginfigm
  drawm("^30,Ph,2'5:N,3:/NH2,4:/COOH,"&
         "2:red,5:blue,3=green")
endfigm
\end{verbatim}
%-----------------------------------------
\begin{mplibcode}
beginfigm
  #(50mm,18mm) |<(8mm)
  drawm("^30,Ph,2'5:N,3:/NH2,4:/COOH,2:red,5:blue,3=green")
endfigm
\end{mplibcode}
%-----------------------------------------------------------------------------
\subsubsection{Change font}
\index{atomfont}%
\begin{verbatim}
beginfigm
  %----------------
  atomfont:="phvro8g";
  %----------------
  drawm("^30,Ph,2'5:N,3:/NH2,4:/COOH")
endfigm
\end{verbatim}
%-----------------------------------------
\begin{mplibcode}
beginfigm
  #(50mm,18mm) |<(8mm)
  atomfont:="phvro8g";
  drawm("^30,Ph,2'5:N,3:/NH2,4:/COOH")
endfigm
\end{mplibcode}
%%%%%%%%%%%%%%%%%%%%%%%%%%%%%%%%%%%%%%%%%%%%%%%%%%%%%%%%%%%%%%%%%%%%%%%%%%%%%%%
\section{Option parameter}
%------------------------------------------------------------------------------
\subsection{Angle parameter}
\index{mangle}%
\begin{verbatim}
mangle=0   ** default

@(0.2,0.5) drawm("Ph")
mangle:=30;
@(0.8,0.5) drawm("Ph")
\end{verbatim}
%-----------------------------------------
%% mangle
\begin{mplibcode}
beginfigm
  #(50mm,15mm) |=(6mm)
  mangle:=0;
  @(0.2,0.5) drawm("Ph")
  ++(drawarrow((A1 shifted (aw,0)) rotated A1ang..A1);)
  mangle:=30;
  @(0.8,0.5) drawm("Ph")
  ++(drawarrow((A1 shifted (aw,0)) rotated A1ang..A1);)
endfigm
\end{mplibcode}
%------------------------------------------------------------------------------
\subsection{Size/Ratio parameter}
%-----------------------------------------------------------------------------
\subsubsection{Bond length [ \textbar={ }]}
\index{blength}%
\index{\textbar=()}%
\begin{verbatim}
|=(n) : abbreviated form of blength:=n;

(fit to figure size)
blength=0   ** default
\end{verbatim}
%-----------------------------------------
\begin{mplibcode}
%%% blength=0
beginfigm
  sw_frame:=Bothside;
  #(40mm,15mm);
  drawm("^30,Ph")
  **(pickup pencircle scaled 0.2pt; 
      for i=0 upto w/mm: draw (i*mm,0)--(i*mm,-.5mm); endfor
      for i=0 upto h/mm: draw (0,i*mm)--(-.5mm,i*mm); endfor
      for i=0 upto w/cm: draw (i*cm,0)--(i*cm,-.8mm); endfor
      for i=0 upto h/cm: draw (0,i*cm)--(-.8mm,i*cm); endfor
  )
endfigm
\end{mplibcode}
%-----------------------------------------
\begin{verbatim}
(ratio bond/figure width)
blength=0.1  ** (0<blength<=1)
blength=60mm(width)*0.1=6mm
\end{verbatim}
%---------------------------------------------
\begin{mplibcode}
%% 0<blength=<1
beginfigm
  sw_frame:=Bothside;
  #(40mm,15mm) |=(0.1)
  drawm("^30,Ph")
  **(pickup pencircle scaled 0.2pt; 
      for i=0 upto w/mm: draw (i*mm,0)--(i*mm,-.5mm); endfor
      for i=0 upto h/mm: draw (0,i*mm)--(-.5mm,i*mm); endfor
      for i=0 upto w/cm: draw (i*cm,0)--(i*cm,-.8mm); endfor
      for i=0 upto h/cm: draw (0,i*cm)--(-.8mm,i*cm); endfor
  )
endfigm
\end{mplibcode}
%-----------------------------------------
\begin{verbatim}
(bond length)
blength=9mm
** (blength>1) ignore msize(w,h)
\end{verbatim}
%---------------------------------------------
\begin{mplibcode}
%%%% blength>1
beginfigm
  sw_frame:=Bothside;
  #(40mm,15mm) |=(8mm)
  drawm("^30,Ph")
  **(pickup pencircle scaled 0.2pt; 
      for i=0 upto w/mm: draw (i*mm,0)--(i*mm,-.5mm); endfor
      for i=0 upto h/mm: draw (0,i*mm)--(-.5mm,i*mm); endfor
      for i=0 upto w/cm: draw (i*cm,0)--(i*cm,-.8mm); endfor
      for i=0 upto h/cm: draw (0,i*cm)--(-.8mm,i*cm); endfor
  )
endfigm
\end{mplibcode}
%------------------------------------------------------------------------------
\subsubsection{Molecular size}
\index{msize}%
\index{\#\#}%
\begin{verbatim}
##(1,1) : msize=(1,1)  ** default
## p : abbreviated form of msize:=p;
\end{verbatim}
%---------------------------------------------
%% msize=(1)
\begin{mplibcode}
beginfigm
  sw_frame:=Bothside+Mol;
  #(40mm,15mm) ##(1,1)
  drawm("^30,Ph")
  **(pickup pencircle scaled 0.2pt; 
      for i=0 upto w/mm: draw (i*mm,0)--(i*mm,-.5mm); endfor
      for i=0 upto h/mm: draw (0,i*mm)--(-.5mm,i*mm); endfor
      for i=0 upto w/cm: draw (i*cm,0)--(i*cm,-.8mm); endfor
      for i=0 upto h/cm: draw (0,i*cm)--(-.8mm,i*cm); endfor
  )
endfigm
\end{mplibcode}
%-----------------------------------------------------------
\begin{verbatim}
##(0.25,1) : msize=(0.25,1)
** msize=(40mm-4mm)*0.25=9mm
\end{verbatim}
%---------------------------------------------
%% msize=(0.25,1)
\begin{mplibcode}
beginfigm
  sw_frame:=Bothside+Mol;
  #(40mm,15mm) ##(0.25,1)
  drawm("^30,Ph")
  **(pickup pencircle scaled 0.2pt; 
      for i=0 upto w/mm: draw (i*mm,0)--(i*mm,-.5mm); endfor
      for i=0 upto h/mm: draw (0,i*mm)--(-.5mm,i*mm); endfor
      for i=0 upto w/cm: draw (i*cm,0)--(i*cm,-.8mm); endfor
      for i=0 upto h/cm: draw (0,i*cm)--(-.8mm,i*cm); endfor
  )
endfigm
\end{mplibcode}
%-----------------------------------------------------------
\begin{verbatim}
##(11mm,11mm) : msize=(11mm,11mm)
\end{verbatim}
%---------------------------------------------
%% msize=(11mm,11mm)
\begin{mplibcode}
beginfigm
  sw_frame:=Bothside+Mol;
  #(40mm,15mm) ##(11mm,11mm)
  drawm("^30,Ph")
  **(pickup pencircle scaled 0.2pt; 
      for i=0 upto w/mm: draw (i*mm,0)--(i*mm,-.5mm); endfor
      for i=0 upto h/mm: draw (0,i*mm)--(-.5mm,i*mm); endfor
      for i=0 upto w/cm: draw (i*cm,0)--(i*cm,-.8mm); endfor
      for i=0 upto h/cm: draw (0,i*cm)--(-.8mm,i*cm); endfor
  )
endfigm
\end{mplibcode}
%------------------------------------------------------------------------------
\subsubsection{Molecular position}
\index{mposition}%
\begin{verbatim}
@(0.5,0.5) : mposition=(0.5,0.5)  **default
\end{verbatim}
%---------------------------------------------
%% mposition
\begin{mplibcode}
beginfigm
  sw_frame:=Bothside+Mol;
  #(40mm,15mm) ##(1,0.8)
  @(0.5,0.5)
  drawm("^30,Ph")
endfigm
\end{mplibcode}
%--------------------------------------------------------------
\begin{verbatim}
@(1,0) : mposition=(1,0)
\end{verbatim}
%---------------------------------------------
%% mposition
\begin{mplibcode}
beginfigm
  sw_frame:=Bothside+Mol;
  #(40mm,15mm) ##(1,0.8)
  @(1,0)
  drawm("^30,Ph")
endfigm
\end{mplibcode}
%--------------------------------------------------------------
\begin{verbatim}
@(10mm,4mm) : mposition=(10mm,4mm)
\end{verbatim}
%---------------------------------------------
\begin{mplibcode}
beginfigm
  sw_frame:=sw_frame+Mol;
  #(40mm,15mm) ##(1,0.8) @(10mm,4mm)
  drawm("^30,Ph")
  **(drawdot p1 withpen pencircle scaled 3pt;
      pickup pencircle scaled 0.2pt; 
      for i=0 upto w/mm: draw (i*mm,0)--(i*mm,-.5mm); endfor
      for i=0 upto h/mm: draw (0,i*mm)--(-.5mm,i*mm); endfor
      for i=0 upto w/cm: draw (i*cm,0)--(i*cm,-.8mm); endfor
      for i=0 upto h/cm: draw (0,i*cm)--(-.8mm,i*cm); endfor
  )
endfigm
\end{mplibcode}
%------------------------------------------------------------------------------
\subsection{Size parameter}
%-----------------------------------------------------------------------------
\subsubsection{Figure size [ \#{ }]}
\index{fsize}%
\index{\#()}%
\begin{verbatim}
fsize=(figure width,figure height)
** default: (30mm,20mm)
# p : abbreviated form of fsize:=p;

#(40mm,15mm) : fsize=(40mm,15mm)
\end{verbatim}
%---------------------------------------------
%% Figure size
\begin{mplibcode}
beginfigm
  sw_frame:=Outside;
  #(40mm,12mm)
  drawm("^30,Ph")
  **(pickup pencircle scaled 0.2pt; 
      for i=0 upto w/mm: draw (i*mm,0)--(i*mm,-.5mm); endfor
      for i=0 upto h/mm: draw (0,i*mm)--(-.5mm,i*mm); endfor
      for i=0 upto w/cm: draw (i*cm,0)--(i*cm,-.8mm); endfor
      for i=0 upto h/cm: draw (0,i*cm)--(-.8mm,i*cm); endfor
  )
endfigm
\end{mplibcode}
%-----------------------------------------------------------------------------
\subsubsection{Figure margin [ \#@{ }]}
\index{fmargin}%
\index{\#"@()}%
\begin{verbatim}
fmargin=(margin left rigth,top bottom)
** default: (0.4mm,0.4mm)
#@ p : abbreviated form of fmargin:=p;

#@(10mm,2mm) : fmargin=(10mm,2mm)
\end{verbatim}
%---------------------------------------------
%% fmargin
\begin{mplibcode}
beginfigm
  sw_frame:=Bothside+Mol;
  #(40mm,12mm) #@(10mm,2mm)
  drawm("^30,Ph")
  **(pickup pencircle scaled 0.2pt; 
      for i=0 upto w/mm: draw (i*mm,0)--(i*mm,-.5mm); endfor
      for i=0 upto h/mm: draw (0,i*mm)--(-.5mm,i*mm); endfor
      for i=0 upto w/cm: draw (i*cm,0)--(i*cm,-.8mm); endfor
      for i=0 upto h/cm: draw (0,i*cm)--(-.8mm,i*cm); endfor
      drawdot p0 withpen pencircle scaled 3pt;
  )
endfigm
\end{mplibcode}
%-----------------------------------------------------------------------------
\subsubsection{Offset thickness of bond}
\index{offset\_thickness}%
\begin{verbatim}
default: offset_thickness=0.25pt
\end{verbatim}
%---------------------------------------------
%%offset_thickness 
\begin{mplibcode}
beginfigm
  #(60mm,12mm);
  offset_thickness:=0.0pt; @(0.1,0.5) drawm("^30,Ph")
  offset_thickness:=0.25pt; @(0.55,0.5) drawm("^30,Ph")
  offset_thickness:=0.5pt; @(1,0.5) drawm("^30,Ph")
  **(defaultscale:=0.6; labeloffset:=1bp;
    label.urt("0.0pt",(0,1bp));
    label.urt("0.25pt",(0.36w,1bp));
    label.urt("0.5pt",(0.7w,1bp));
  )
endfigm
\end{mplibcode}
%-----------------------------------------------------------------------------
\subsubsection{Offset of double bond gap}
\index{offset\_bond\_gap}%
\begin{verbatim}
default: offset_bond_gap=0.3pt
\end{verbatim}
%---------------------------------------------
\begin{mplibcode}
beginfigm
  #(60mm,12mm);
  offset_bond_gap:=0.0pt; @(0.1, 0.5) drawm("^30,Ph")
  offset_bond_gap:=0.3pt; @(0.55,0.5) drawm("^30,Ph")  %<<== default
  offset_bond_gap:=1.0pt; @(1,   0.5) drawm("^30,Ph")
  **(defaultscale:=0.6; labeloffset:=1bp;
    label.urt("0.0pt",(0,1bp));
    label.urt("0.3pt",(0.36w,1bp));
    label.urt("1.0pt",(0.7w,1bp));
  )
endfigm
\end{mplibcode}
%-----------------------------------------------------------------------------
\subsubsection{Offset of atom width}
\index{offset\_atom}%
\begin{verbatim}
default: offset_atom=0.8pt
\end{verbatim}
%---------------------------------------------
\begin{mplibcode}
beginfigm
  #(60mm,12mm);
  offset_atom:=0.0pt; @(0.1, .5) drawm("^30,?6,3:O")
  offset_atom:=0.8pt; @(.55, .5) drawm("^30,?6,3:O")   %<<== default
  offset_atom:=2.0pt; @(1,   .5) drawm("^30,?6,3:O")
  **(defaultscale:=0.6; labeloffset:=1bp;
    label.urt("0.0pt",(0,1bp));
    label.urt("0.8pt",(0.36w,1bp));
    label.urt("2.0pt",(0.7w,1bp));
  )
endfigm
\end{mplibcode}
%-----------------------------------------------------------------------------
\subsubsection{Offset of wedge width}
\index{offset\_wedge}%
\begin{verbatim}
default: offset_wedge=0.2pt
\end{verbatim}
%---------------------------------------------
\begin{mplibcode}
beginfigm
  #(60mm,12mm);
  offset_wedge:=0.0pt; @(0.1,0.5) drawm("^30,?6,5:?w")
  offset_wedge:=0.2pt; @(0.55,0.5) drawm("^30,?6,5:?w")  %<<== default
  offset_wedge:=0.5pt; @(1, 0.5) drawm("^30,?6,5:?w")
  **(defaultscale:=0.6; labeloffset:=1bp;
    label.urt("0.0pt",(0,1bp));
    label.urt("0.2pt",(0.36w,1bp));
    label.urt("0.5pt",(0.7w,1bp));
  )
endfigm
\end{mplibcode}
%-----------------------------------------------------------------------------
\subsubsection{Max bond length [ \textbar\textless{ }]}
\index{max\_blength}%
\index{\textbar\textless()}%
\begin{verbatim}
|<(n): abbreviated form of max_blength:=n;
default: max_blength=10mm
\end{verbatim}
%---------------------------------------------
\begin{mplibcode}
beginfigm
  #(60mm,12mm)
  sw_frame:=sw_frame+Mol;
  |<(5mm)  @(0, .5) drawm("^45,?4")
  |<(8mm)  @(.4,.5) drawm("^45,?4")
  |<(10mm) @(1, .5) drawm("^45,?4")   %<<== default
  **(defaultscale:=0.6; labeloffset:=1bp;
    label("5mm", (0.075w,0.5h));
    label("8mm", (0.42w,0.5h));
    label("10mm",(0.88w,0.5h));
  )
endfigm
\end{mplibcode}
%-----------------------------------------------------------------------------
\subsection{Ratio parameter}
%-----------------------------------------------------------------------------
\subsubsection{Thickness/bond length}
\index{ratio\_thickness\_bond}%
\begin{verbatim}
default: ratio_thickness_bond=0.012
\end{verbatim}
%---------------------------------------------
\begin{mplibcode}
beginfigm
  #(60mm,12mm);
  ratio_thickness_bond:=0.005;  @(0.1,0.5) drawm("^30,Ph")
  ratio_thickness_bond:=0.012;  @(.55,0.5) drawm("^30,Ph")  %<<== default
  ratio_thickness_bond:=0.03;   @(1,  0.5) drawm("^30,Ph")
  **(defaultscale:=0.6; labeloffset:=1bp;
    label.urt("0.005",(0,1bp));
    label.urt("0.015",(0.36w,1bp));
    label.urt("0.030",(0.7w,1bp));
  )
endfigm
\end{mplibcode}
%-----------------------------------------------------------------------------
\subsubsection{Atom/bond length}
\index{ratio\_atom\_bond}%
\begin{verbatim}
default: ratio_atom_bond= 0.36
\end{verbatim}
%---------------------------------------------
\begin{mplibcode}
beginfigm
  #(60mm,12mm);
  ratio_atom_bond:=0.25;  @(0.1,0.5) drawm("^30,?6,3:O")
  ratio_atom_bond:=0.33;  @(.55,0.5) drawm("^30,?6,3:O")  %<<== default
  ratio_atom_bond:=0.45;  @(1,  0.5) drawm("^30,?6,3:O")
  **(defaultscale:=0.6; labeloffset:=1bp;
    label.urt("0.25",(0,1bp));
    label.urt("0.33",(0.36w,1bp));
    label.urt("0.45",(0.7w,1bp));
  )
endfigm
\end{mplibcode}
%-----------------------------------------------------------------------------
\subsubsection{Char thickness/Atom width}
\index{ratio\_char\_atom}%
\begin{verbatim}
default: ratio_char_atom=0.12
\end{verbatim}
%---------------------------------------------
\begin{mplibcode}
beginfigm
  #(60mm,12mm)
  ratio_char_atom:=0.1;  @(0, 0.5) drawm("^30,?6,6:O,3:NH")
  ratio_char_atom:=0.12; @(.5,0.5) drawm("^30,?6,6:O,3:NH")  %<<== default
  ratio_char_atom:=0.15; @( 1,0.5) drawm("^30,?6,6:O,3:NH")
  **(defaultscale:=0.6; labeloffset:=1bp;
    label.urt("0.1", (0,1bp));
    label.urt("0.12",(0.36w,1bp));
    label.urt("0.15",(0.7w,1bp));
  )
endfigm
\end{mplibcode}
%-----------------------------------------------------------------------------
\subsubsection{Bond gap/bond length}
\index{ratio\_bondgap\_bond}%
\begin{verbatim}
default: ratio_bondgap_bond= 0.15
\end{verbatim}
%---------------------------------------------
\begin{mplibcode}
beginfigm
  #(60mm,12mm);
  ratio_bondgap_bond:=0.10; @(0.1,0.5) drawm("^30,Ph")
  ratio_bondgap_bond:=0.15; @(.55,0.5) drawm("^30,Ph")  %<<== default
  ratio_bondgap_bond:=0.20; @(1  ,0.5) drawm("^30,Ph")
  **(defaultscale:=0.6; labeloffset:=1bp;
    label.urt("0.10",(0,1bp));
    label.urt("0.15",(0.36w,1bp));
    label.urt("0.20",(0.7w,1bp));
  )
endfigm
\end{mplibcode}
%-----------------------------------------------------------------------------
\subsubsection{Wedge/bond length}
\index{ratio\_wedge\_bond}%
\begin{verbatim}
default: ratio_wedge_bond=0.12
\end{verbatim}
%------------------------------------------------------
\begin{mplibcode}
beginfigm
  #(70mm,12mm);
  ratio_wedge_bond:=0.1;  @(0.05,.5) drawm("?6,4:?w")
  ratio_wedge_bond:=0.12; @(.55, .5) drawm("?6,4:?w")  %<<== default
  ratio_wedge_bond:=0.2;  @(1  , .5) drawm("?6,4:?w")
  **(defaultscale:=0.6; labeloffset:=1bp;
    label.urt("0.10",(0,1bp));
    label.urt("0.12",(0.36w,1bp));
    label.urt("0.20",(0.7w,1bp));
  )
endfigm
\end{mplibcode}
%-----------------------------------------------------------------------------
\subsubsection{Figure atom gap/atom length}
\index{ratio\_atomgap\_atom}%
\index{readm}%
\begin{verbatim}
default:  ratio_atomgap_atom= 0.050
\end{verbatim}
%------------------------------------------------------
\begin{mplibcode}
beginfigm
  readm("^30,!2`0.5,2:O")
  #(70mm,12mm)
  sw_frame:=Atom;
  ratio_atomgap_atom:=0.00;  @(0, .5) drawm(mc)
  ratio_atomgap_atom:=0.050; @(.5,.5) drawm(mc)  %<<== default
  ratio_atomgap_atom:=0.12;  @(1, .5) drawm(mc)
  **(defaultscale:=0.6; labeloffset:=1bp;
    label.urt("0.00",(0.05w,3bp));
    label.urt("0.05",(0.45w,3bp));
    label.urt("0.12",(0.85w,3bp));
  )
endfigm
\end{mplibcode}
%-----------------------------------------------------------------------------
\subsubsection{Chain/ring length}
\index{ratio\_chain\_ring}%
\begin{verbatim}
default:  ratio_chain_ring= 0.66
\end{verbatim}
%------------------------------------------------------
\begin{mplibcode}
beginfigm
  #(70mm,12mm);
  ratio_chain_ring:= 0.4;  @(0.05,.5) drawm("^30,?6,4:/!")
  ratio_chain_ring:= 0.66; @(.45, .5) drawm("^30,?6,4:/!")   %<<== default
  ratio_chain_ring:= 1;    @(1,   .5) drawm("^30,?6,4:/!")
  **(defaultscale:=0.6; labeloffset:=1bp;
    label.urt("0.40",(0,1bp));
    label.urt("0.66",(0.3w,1bp));
    label.urt("1.0" ,(0.62w,1bp));
  )
endfigm
\end{mplibcode}
%-----------------------------------------------------------------------------
\subsubsection{Hash gap/bond length}
\index{ratio\_hashgap\_bond}%
\begin{verbatim}
default:  ratio_hashgap_bond=0.12
\end{verbatim}
%------------------------------------------------------
\begin{mplibcode}
beginfigm
  readm("^30,!2,2:?z`1.5")
  #(70mm,15mm)
  ratio_hashgap_bond:=0.06; @(0.08,.5) drawm(mc)
  ratio_hashgap_bond:=0.12; @( .55,.5) drawm(mc) %<<== default
  ratio_hashgap_bond:=0.20; @(1,   .5) drawm(mc)
  **(defaultscale:=0.6; labeloffset:=1bp;
    label.urt("0.06",(0,1bp));
    label.urt("0.12",(0.4w,1bp));
    label.urt("0.20",(0.77w,1bp));
  )
endfigm
\end{mplibcode}
%-----------------------------------------------------------------------------
%%%%\newpage
%-----------------------------------------------------------------------------
\subsection{Drawing mode}
%-----------------------------------------------------------------------------
\subsubsection{Numbering atom,bond}
\index{sw\_numbering}%
\index{Atom}%
\begin{verbatim}
sw_numbering:=Atom;
\end{verbatim}
%------------------------------------------------------
\begin{mplibcode}
beginfigm
  #(60mm,7mm);
  ratio_chain_ring:=1;
  sw_numbering:=Atom; drawm("^-30,!9")
endfigm
\end{mplibcode}
%-----------------------------------------------------------------------------
\index{Bond}%
\begin{verbatim}
sw_numbering:=Bond;
\end{verbatim}
%------------------------------------------------------
\begin{mplibcode}
beginfigm
  #(60mm,7mm);
  ratio_chain_ring:=1;
  sw_numbering:=Bond; drawm("^-30,!9")
endfigm
\end{mplibcode}
%-----------------------------------------------------------------------------
\begin{verbatim}
sw_numbering:=Atom+Bond;
\end{verbatim}
%------------------------------------------------------
\begin{mplibcode}
beginfigm
  #(60mm,7mm);
  ratio_chain_ring:=1;
  sw_numbering:=Atom+Bond; drawm("^-30,!9")
endfigm
\end{mplibcode}
%------------------------------------------------------
\begin{verbatim}
sw_numbering:=Atom+Bond;
\end{verbatim}
%------------------------------------------------------
\begin{mplibcode}
beginfigm
  #(60mm,12mm);
  ratio_chain_ring:=1;
  sw_numbering:=Atom+Bond; drawm("^-30,!9,3'5'7'9:?w`.66")
endfigm
\end{mplibcode}
%------------------------------------------------------
\begin{verbatim}
sw_numbering:=Atom+Bond+All;
\end{verbatim}
%------------------------------------------------------
\begin{mplibcode}
beginfigm
  #(60mm,12mm);
  ratio_chain_ring:=1;
  sw_numbering:=Atom+Bond+All;
  drawm("^-30,!9,3'5'7'9:?w`.66")
endfigm
\end{mplibcode}
%-----------------------------------------------------------------------------
\subsubsection{Trimming mode}
\index{sw\_trimming}%
\begin{verbatim}
sw_trimming:=0;  ** default
##(1,0.7)
@(0.2,0.3) drawm("Ph")
@(0.8,0.7) drawm("Ph")
\end{verbatim}
%------------------------------------------------------
\begin{mplibcode}
beginfigm
  #(60mm,15mm) ##(1,.7)
  sw_frame:=Bothside+Mol;
  @(.2,.3) drawm("Ph")
  @(.8,.7) drawm("Ph")
endfigm
\end{mplibcode}
%------------------------------------------------------
\begin{verbatim}
sw_trimming:=1;
@(0.2,0.3) drawm("Ph")
@(0.8,0.7) drawm("Ph")
\end{verbatim}
%------------------------------------------------------
\begin{mplibcode}
beginfigm
  #(60mm,15mm) ##(1,.7)
  sw_frame:=Bothside+Mol;
  sw_trimming:=1;
  @(.2,.3) drawm("Ph")
  @(.8,.7) drawm("Ph")
endfigm
\end{mplibcode}
%-----------------------------------------------------------------------------
\subsubsection{Omit group}
\index{Group}%
\begin{verbatim}
** default: sw_omit=Group
\end{verbatim}
%------------------------------------------------------
\begin{mplibcode}
beginfigm
  readm("^30,Ph,4:/Cl,3:/F");
  #(60mm,12mm)
  @(.15, .5) putm
  sw_omit:=Group;
  @(.85, .5) putm
endfigm
\end{mplibcode}
%-----------------------------------------------------------------------------
\subsubsection{Omit bond type}
\index{Bond}%
\begin{verbatim}
** default: sw_omit=Bond
\end{verbatim}
%------------------------------------------------------
\begin{mplibcode}
beginfigm
  readm("^30,Ph,4:/Cl,3:/F");
  #(60mm,12mm)
  @(.15, .5) putm
  sw_omit:=Bond;
  @(.85, .5) putm
endfigm
\end{mplibcode}
%-----------------------------------------------------------------------------
\subsection{Frame}
%-----------------------------------------------------------------------------
\subsubsection{Figure frame}
\index{sw\_frame}%
\index{Bothside}%
\index{Inside}%
\index{Outside}%
\begin{verbatim}
** default:sw_frame=0
(Draw figure frame)
fmargin:=(5mm,2mm);
sw_frame=Outside
\end{verbatim}
%------------------------------------------------------
\begin{mplibcode}
beginfigm
  #(30mm,10mm) #@(5mm,1.5mm)
  sw_frame:=Outside;
  drawm("^30,Ph")
endfigm
\end{mplibcode}
%------------------------------------------------------
\begin{verbatim}
(Frame inside margin)
sw_frame=Inside
\end{verbatim}
%------------------------------------------------------
\begin{mplibcode}
beginfigm
  #(30mm,10mm) #@(5mm,1.5mm)
  sw_frame:=Inside;
  drawm("^30,Ph")
endfigm
\end{mplibcode}
%------------------------------------------------------
\begin{verbatim}
(Draw both frame)
sw_frame=Bothside=Inside+Outside
\end{verbatim}
%------------------------------------------------------
\begin{mplibcode}
beginfigm
  #(30mm,10mm) #@(5mm,1.5mm)
  sw_frame:=Bothside;
  drawm("^30,Ph")
endfigm
\end{mplibcode}
%-----------------------------------------------------------------------------
\subsubsection{Molecular frame}
\index{Mol}%
\begin{verbatim}
sw_frame=Mol
** default:sw_frame=0
\end{verbatim}
%------------------------------------------------------
\begin{mplibcode}
beginfigm
  sw_frame:=Outside;
  #(40mm,11mm) ##(1,1)
  sw_frame:=sw_frame+Mol;
  drawm("^30,Ph")
endfigm
\end{mplibcode}
%-----------------------------------------------------------------------------
\subsubsection{Atom frame}
\index{Atom}%
\begin{verbatim}
sw_frame=Atom
** default: sw_frame=0

drawm("^30,COOH,!,COOH")
\end{verbatim}
%---------------------------------------------
\begin{mplibcode}
beginfigm
  #(60mm,10mm);
  sw_frame:=Atom;
  drawm("^30,COOH,!,COOH")
endfigm
\end{mplibcode}
%-----------------------------------------------------------------------------
\subsection{Parameter setting}
\subsubsection{Local parameter setting}
\index{beginfigm}%
\index{endfigm}%
\begin{verbatim}
beginfigm
  drawm("Ph")
endfigm
beginfigm
  %--------------------------
  ratio_thickness_bond:=0.05;
  %--------------------------
  drawm("Ph")
endfigm
beginfigm 
  drawm("Ph")
endfigm
\end{verbatim}
%------------------------------------------------------
\quad
\begin{mplibcode}
fsize:=(15mm,12mm);
beginfigm
  drawm("Ph")
endfigm
beginfigm
  ratio_thickness_bond:=0.05;
  drawm("Ph")
endfigm
beginfigm
  drawm("Ph")
endfigm
\end{mplibcode}
%-----------------------------------------------------------------------------
\newpage
\subsubsection{Global parameter setting}
\begin{verbatim}
beginfigm
  drawm("Ph")
endfigm
%--------------------------
ratio_thickness_bond:=0.05;
%--------------------------
beginfigm
  drawm("Ph")
endfigm
beginfigm
  drawm("Ph")
endfigm
beginfigm
  drawm("Ph")
endfigm
\end{verbatim}
%----------------------------------------------------
\quad
\begin{mplibcode}
save_ratio:=ratio_thickness_bond;
fsize:=(15mm,12mm);
beginfigm
  drawm("Ph")
endfigm
ratio_thickness_bond:=0.05;
beginfigm
  drawm("Ph")
endfigm
beginfigm
  drawm("Ph")
endfigm
beginfigm
  drawm("Ph")
endfigm
ratio_thickness_bond:=save_ratio;
\end{mplibcode}
%%%%%%%%%%%%%%%%%%%%%%%%%%%%%%%%%%%%%%%%%%%%%%%%%%%%%%%%%%%%%%%%%%%%%%%%%%%%%%%
\section{Command of mcf2graph}
%-----------------------------------------------------------------------------
\subsection{drawm}
\index{drawm}%
\begin{verbatim}
(Draw molecule)

msize=(a,b)        **default (1,1)
mposition=(c,d)    **default (0.5,0.5)

a: ratio molecular width/figure width
b: ratio molecular hight/figure hight
c: x axis position
d: y axis position

  drawm("^30,Ph,3:/F,4:/Cl")
\end{verbatim}
%------------------------------------------------
\begin{mplibcode}
beginfigm
  #(40mm,15mm) #@(0.5mm,0.5mm)
  sw_frame:=Outside+Mol;
  drawm("^30,Ph,4:/Cl,3:/F")
endfigm
\end{mplibcode}
%-----------------------------------------------------------------------------
\subsection{readm}
\index{readm()}%
\begin{verbatim}
readm(string1,string2, ...);
** string = mcf code

(example)
readm("^30,Ph,1'2'6:/O!,-4'-5=?7,",
   " -1'-4'-6=dl,-2:?O,-3:/O!,   ",
   " 9:,!,NH,!,?O,!              ");
putm
\end{verbatim}

\begin{mplibcode}
beginfigm
#(40mm,20mm);
readm("^30,Ph,1'2'6:/O!,-4'-5=?7,",
   " -1'-4'-6=dl,-2:?O,-3:/O!,   ",
   " 9:,!,NH,!,?O,!              ");
putm
endfigm
\end{mplibcode}
%-----------------------------------------------------------------------------
\subsection{getm}
\index{getm}%
\index{ucount}%
\begin{verbatim}
getm(number)
** number = numeric
** ucount = molecular data unit count

for i=1 upto ucount:
  beginfigm
    getm(i)        % get data unit no=i
    putm           % put figure
  endfigm
endfor

getm("name"): "name"=string

(example)
 beginfigm
   getm("Adenine")
   putm
 endfigm
\end{verbatim}
%-----------------------------------------------------------------------------
\subsection{putm}
\index{putm}%
\begin{verbatim}
putm : put figure

  if op_row>=1: scantokens(op) fi
  if mc_row>=1: drawm(mc) fi
  if ad_row>=1: add(scantokens(ad)) fi
  if ex_row>=1: ext(scantokens(ex)) fi
\end{verbatim}
%-------------------------------------------------------------------------------
\subsection{allm}
\index{allm}%
\begin{verbatim}
allm : put all loaded molecule

  for i=1 upto ucount:
    beginfigm getm(i) putm endfigm
  endfor
\end{verbatim}
%===============================================================================
\newpage
\subsection{add [ ++ ]}
\index{add}%
\index{++}%
\index{plus}%
\index{minus}%
\index{lonepair}%
\index{lonepairdiam}%
\index{lonepairspace}%
\index{circlediam}%
\index{circlepen}%
\index{w}%
\index{h}%
\index{aw}%
\index{em}%
\index{p0}%
\index{l}%
\index{/*}%
\index{**}%
\index{\textgreater\textgreater}%
\index{An}%
\index{A[]}%
\index{A[]ang}%
\index{A[]up}%
\index{A[]left}%
\index{A[]right}%
\index{A[]down}%
\index{Bn}%
\index{B[]}%
\index{B[]s}%
\index{B[]m}%
\index{B[]e}%
\index{B[]ang}%
\index{B[]up}%
\index{B[]left}%
\index{B[]right}%
\index{B[]down}%
\index{defaultscale}%
\index{labeloffset}%
\begin{verbatim}
(Add label to molecule)

++(): add()
w:       molecular width
h:       molecular height
aw:      atom font size
em:      label font size
p0:      origin of molecular structure
l:       bond length

An:      atom number
A[m]:    atom position
A[m]ang: branch angle of A[m]
A[m]up:    dir A[m]ang
A[m]left:  dir A[m]ang+90
A[m]right: dir A[m]ang-90
A[m]down:  dir A[m]ang+180

Bn:      bond number
B[m]:    bond(path)
B[m]s:   bond start position
B[m]m:   bond middle position
B[m]e:   bond end position
B[m]ang: bond angle
B[m]up:    dir B[m]ang
B[m]left:  dir B[m]ang+90
B[m]right: dir B[m]ang-90
B[m]down:  dir B[m]ang+180

plus  : '+' circled
minus : '-' circled
  circlediam = 0.6aw (default)
  circlepen =  0.2bp (default)

lonepair r: ':' rotated r
  lonepairdiam  = 0.3aw (default)
  lonepairspace = 0.7aw (default)

# : scaled
<< : rotated
a /* b : point b of a

beginfigm
 #(70mm,40mm) ##(.91,.9) |<(10mm)
 sw_frame:=sw_frame+Atom+Mol;
 @(.5,.85) drawm("^30,?6,2'5:O")
 ++(
  defaultscale:=.8;
  labeloffset:=.3aw;
  dotlabel.lft("p0",p0);
  dotlabel.rt( "p0+(w,h)",p0+(w,h));
  dotlabel.ulft("A1",A1);
  drawarrow A1..A1+__*l<<A1ang;
  dotlabel.lrt( "B3s",B3s);
  dotlabel.rt("B3m",B3m);
  drawarrow B3m..B3m+__*l<<(B3ang+90);
  dotlabel.ulft("A6",A6);
  drawarrow A1{A1down}..A6;
  dotlabel.urt( "B3e",B3e);
  label.rt(  "An="&decimal(An)&
           "  Bn="&decimal(Bn)&
           "  aw="&decimal(aw)&
           "  em="&decimal(em),
           p0+(-9em,-1.5em));
  label.rt(  "w="&decimal(w)&
    "  h="&substring (0,6)of decimal(h)&
    "  l="&substring (0,6)of decimal(l),
           p0+(-9em,-3em));
 )
endfigm
\end{verbatim}
%------------------------------------------------------------------------------
\begin{mplibcode}
beginfigm
 #(70mm,40mm) ##(.91,.9) |<(10mm)
 sw_frame:=sw_frame+Atom+Mol;
 @(.5,.85) drawm("^30,?6,2'5:O")
 ++(
  defaultscale:=.8;
  labeloffset:=.3aw;
  dotlabel.lft("p0",p0);
  dotlabel.rt( "p0+(w,h)",p0+(w,h));
  dotlabel.ulft("A1",A1);
  drawarrow A1..A1+__*l<<A1ang;
  dotlabel.lrt( "B3s",B3s);
  dotlabel.rt("B3m",B3m);
  drawarrow B3m..B3m+__*l<<(B3ang+90);
  dotlabel.ulft("A6",A6);
  drawarrow A1{A1down}..A6;
  dotlabel.urt( "B3e",B3e);
  label.rt(  "An="&decimal(An)&
           "  Bn="&decimal(Bn)&
           "  aw="&decimal(aw)&
           "  em="&decimal(em),
           p0+(-9em,-1.5em));
  label.rt(  "w="&substring (0,6)of decimal(w)&
           "  h="&substring (0,6)of decimal(h)&
           "  l="&substring (0,6)of decimal(l),
           p0+(-9em,-3em));
 )
endfigm
\end{mplibcode}
%------------------------------------------------------------------------------
\begin{verbatim}
beginfigm
 #(60mm,20mm) ##(1,0.85)
 %---------------------------------------
 @(0,0) drawm("^30,Ph,3=dl,4:/NH2")
 %---------------------------------------
 ++(
  labeloffset:=.7aw;
  label.top(lone_pair 90,A7);
  drawarrow 
    (A7+up#1.2aw){A7left}
     ..{B7right}B7/*0.3;
  drawarrow
    B3m..A3+B2up#1.5aw..{A3down}A3;
 )
 %---------------------------------------
 @(1,0) drawm("^30,?6,1'5=dl,4://NH2")
 %---------------------------------------
 ++(
  labeloffset:=.7aw;
  label.top(plus,A7);
  label.urt(minus,A3);
  label(lonepair A3ang,A3+A3up#.7aw);
 )
 %---------------------------------------
 **(drawdblarrow (.4w,.4h)..(.55w,.4h);)
 %---------------------------------------
endfigm
\end{verbatim}
%-------------------------------------------------------
\begin{mplibcode}
beginfigm
  #(70mm,20mm) ##(1,0.85)
  %-------------------------------------------
  @(0,0) drawm("^30,Ph,3=dl,4:/NH2")
  %-------------------------------------------
  ++(labeloffset:=.7aw;
      label.top(lonepair 90,A7);
      drawarrow (A7+up#1.2aw){A7left}..{B7right}B7/*0.3;
      drawarrow B3m..A3+B2up#1.5aw..{A3down}A3;
  )
  %-------------------------------------------
  @(1,0) drawm("^30,?6,1'5=dl,4://NH2")
  %-------------------------------------------
  ++(labeloffset:=.7aw;
      label.top(plus,A7);
      label.urt(minus,A3);
      label(lonepair A3ang,A3+A3up#.7aw);
  )
  **(drawdblarrow (0.4w,0.4h)..(0.55w,0.4h);)
endfigm
\end{mplibcode}
%-----------------------------------------------------------------------------
\newpage
\subsection{ext [ ** ]}
\index{ext}%
\index{**}%
\index{w0}%
\index{h0}%
\index{aw}%
\index{em}%
\index{n}%
\index{ratio\_thickness\_char}%
\index{defaultscale}%
\begin{verbatim}
(Extra label to figure)
 
**(): ext()
w:    figure width
h:    figure height
w0:   figure width-2xpart(fmargin)
h0:   figure height-2ypart(fmargin)
aw:   atom font size
em:   label font size
p0:   fmargin

n:    molecular number
p[m]: molecular origin position
w[m]: molecular width
h[m]: molecular height

ratio_thickness_char:
pen thickness / char width
%----------------------------------------
beginfigm
 #(70mm,30mm)
 |=(0.065)
 %---------------------------------------
 @(0.1,0.5)
 drawm("^-210,60`1,60`1,60`1,1'3=dl")
 ++(
   defaultscale:=0.6;
   label.bot("Diene",p0+(0.5w,0));
   )
 @(0.4,0.5)
 drawm("^-30,-60`1,1=dm")
 ++(
   defaultscale:=0.6;
   label.bot("Dienophile",p0+(.5w,0));
   )
 @(0.9,0.5) drawm("^30,?6,6=dl")
 %---------------------------------------
 **(
  drawarrow (.52w,.5h)..(.6w,.5h);
  defaultscale:=0.7;
  label("+",(0.25w,0.5h));
  ratio_thickness_char:=0.125;
  label.bot("Diels-Alder Reaction",
            (.5w,h)
           );
 )
 %---------------------------------------
endfigm
\end{verbatim}
%------------------------------------------------------------------------------
\begin{mplibcode}
beginfigm
 #(70mm,30mm)
 |=(0.065)
 %---------------------------------------
 @(0.1,0.5)
 drawm("^-210,60`1,60`1,60`1,1'3=dl")
 ++(
   defaultscale:=0.6;
   label.bot("Diene",p0+(0.5w,0));
 )
 @(0.4,0.5)
   drawm("^-30,-60`1,1=dm")
   ++(defaultscale:=0.6;
   label.bot("Dienophile",p0+(.5w,0));
 )
 @(0.9,0.5) drawm("^30,?6,6=dl")
 %---------------------------------------
 **(
  drawarrow (.52w,.5h)..(.6w,.5h);
  defaultscale:=0.7;
  label("+",(0.25w,0.5h));
  ratio_thickness_char:=0.125;
  label.bot("Diels-Alder Reaction",
            (.5w,h));
 )
 %---------------------------------------
endfigm
\end{mplibcode}
%-----------------------------------------------------------------------------
\subsubsection{Local ext setting}
\index{\_}%
\begin{verbatim}
beginfigm
  EN:="?3";@(0.5,1) drawm("^30,?3")
endfigm
beginfigm
  EN:="?4";@(0.5,1) drawm("?4")
  %-------------------------------
  **(label.top(EN,(0.5w,0));)
  %-------------------------------
endfigm
beginfigm
  EN:="?5";@(0.5,1) drawm("?5")
endfigm
beginfigm
  EN:="?6";@(0.5,1) drawm("?6")
endfigm
\end{verbatim}
%----------------------------------------------------
\quad
\begin{mplibcode}
beginfigm
  EN:="?3";#(12mm,15mm)@(0.5,1) drawm("^30,?3")
endfigm
beginfigm
  EN:="?4";#(12mm,15mm)@(0.5,1) drawm("?4")
  %-------------------------------
  **(label.top(EN,(0.5w,0));)
  %-------------------------------
endfigm
beginfigm
  EN:="?5";#(12mm,15mm)@(0.5,1) drawm("?5")
endfigm
beginfigm
  EN:="?6";#(12mm,15mm)@(0.5,1) drawm("?6")
endfigm
\end{mplibcode}
%-----------------------------------------------------------------------------
\subsubsection{Global ext setting}
\index{ext\_clear}%
\begin{verbatim}
ext_clear: reset global ext()

beginfigm
  EN:="?3";@(0.5,1) drawm("^30,?3")
endfigm
%-------------------------------
ext(label.top(EN,(0.5w,0));)
%-------------------------------
beginfigm
  EN:="?4";@(0.5,1) drawm("?4")
endfigm
beginfigm
  EN:="?5";@(0.5,1) drawm("?5")
endfigm
%---------
ext_clear;
%---------
beginfigm
  EN:="?6";@(0.5,1) drawm("?6")
endfigm
\end{verbatim}
%-------------------------------------
\quad
\begin{mplibcode}
beginfigm
  EN:="?3";#(12mm,15mm)
  @(0.5,1) drawm("^30,?3")
endfigm
%-------------------------------
ext(label.top(EN,(0.5w,0));)
%-------------------------------
beginfigm
  EN:="?4";#(12mm,15mm)
  @(0.5,1) drawm("?4")
endfigm
beginfigm
  EN:="?5";#(12mm,15mm)
  @(0.5,1) drawm("?5")
endfigm
%---------
ext_clear;
%---------
beginfigm
 EN:="?6";#(12mm,15mm)
  @(0.5,1) drawm("?6")
endfigm
\end{mplibcode}
%-----------------------------------------------------------------------------
\newpage
%-----------------------------------------------------------------------------
\onecolumn
\section{Example}
%-----------------------------------------------------------------------------
\subsection{drawm example}
\paragraph{(Luciferin)}
\begin{verbatim}
beginfigm
  #(50mm,15mm)
  drawm("^30,Ph,3=?5,8:,!,?5,9'16=dl,"&
        "9'14:N,7'11:S,1:/OH,-2:*/COOH")
endfigm
\end{verbatim}
%-------------------------------------------
\begin{mplibcode}
loadm("EN<>*");
beginfigm
  #(50mm,15mm)
  drawm("^30,Ph,3=?5,8:,!,?5,9'16=dl,"&
        "9'14:N,7'11:S,1:/OH,-2:*/COOH")
endfigm
\end{mplibcode}
%-----------------------------------------------------------------------------
\subsection{readm example}
%-----------------------------------------------------------------------------
\paragraph{(Colchicine)}
\begin{verbatim}
beginfigm
  readm(
  "^30,Ph,1'2'6:/O!,-4'-5=?7,     ",
  " -1'-4'-6=dl,-2:?O,-3:/O!,     ",
  " 9:/NH!'?O!                    ")
  #(40mm,20mm) putm
endfigm
\end{verbatim}
%---------------------------------------
% Colchicine MW:385.41
\begin{mplibcode}
beginfigm
  readm(
  "^30,Ph,1'2'6:/O!,-4'-5=?7,    ",
  " -1'-4'-6=dl,-2:?O,-3:/O!,    ",
  " 9:/NH!'?O!                   ")
  #(40mm,20mm) putm
endfigm
\end{mplibcode}
%-----------------------------------------------------------------------------
\paragraph{(Maltose)}
\index{arc\_lb}
\index{arc\_br}
\begin{verbatim}
(bond type for glycan)
arc_lb : arc left > bottom
arc_br : arc bottom right

beginfigm
  %"EN:Maltose","MW:342.3",
  readm(
  %--------------------------------------------------------
  "hexose_hp,`.5,1^$270'2^$90'3^$270:/OH,6^$90:/!OH,`,   ",
  "4:,$310~arc_lb`1,O,$50~arc_br`1,^$0,                  ",
  "<,hexose_hp,`.5,2^$90'3^$270'4^$270:/OH,6^$90:/!OH    ")
  %--------------------------------------------------------
  #(50mm,20mm) putm
endfigm
\end{verbatim}
%-------------------------------------------------
\begin{mplibcode}
beginfigm
  %"EN:Maltose","MW:342.3",
  readm(
  %--------------------------------------------------------
  "hexose_hp,`.5,1^$270'2^$90'3^$270:/OH,6^$90:/!OH,`,   ",
  "4:,$310~arc_lb`1,O,$50~arc_br`1,^$0,                  ",
  "<,hexose_hp,`.5,2^$90'3^$270'4^$270:/OH,6^$90:/!OH    ")
  %--------------------------------------------------------
  #(50mm,20mm) putm
endfigm
\end{mplibcode}
%-----------------------------------------------------------------------------
\newpage
\paragraph{(Erythromycin)}
\begin{verbatim}
beginfigm
  EN:=Erythromycin"; MW:="733.93";
  #(60mm,30mm) @(1,0.5)
  readm(
  %-----------------------------------------------------------------------
    "^30,`1,^-120,60,60,60,-60,60,60,-60,60,60,60,-60,60,60,`,:1,       ",
    " 14:O,13:/*!,1'9:?O,                                               ",
    " *2'4'6^-35'8'*10'12^35:?z,                                        ",
    " 6^35'11'12^-35:*/OH,                                              ",
    " $3:,!z,O,30~zb,<,?6`.7,6:O,`.5,5~wf'3^35:?,4:/*OH,3^-35:/*O!,`,   ",
    " $5:,!z^30`1.7,O,!~zb,<,?6`.7,6:O,`.5,5:?z,2:*/OH,3:/*N?!          ")
  %-----------------------------------------------------------------------
  putm
  **(defaultscale:=0.8;
      label.lrt("EN: "&EN,(0,h));
      label.lrt("fm: "&fm,(0,h-5mm));
      label.lrt("mw: "&mw,(0,h-9mm));
      label.lrt("MW: "&MW,(0,h-13mm));
  )
endfigm
\end{verbatim}
%----------------------------------------------------------------------------
\begin{mplibcode}
beginfigm
  EN:="Erythromycin"; MW:="733.93";
  #(60mm,30mm) @(1,0.5)
  readm(
  %-----------------------------------------------------------------------
    "^30,`1,^-120,60,60,60,-60,60,60,-60,60,60,60,-60,60,60,`,:1,       ",
    " 14:O,13:/*!,1'9:?O,                                               ",
    " *2'4'6^-35'8'*10'12^35:?z,                                        ",
    " 6^35'11'12^-35:*/OH,                                              ",
    " $3:,!z,O,30~zb,<,?6`.7,6:O,`.5,5~wf'3^35:?,4:/*OH,3^-35:/*O!,`,   ",
    " $5:,30~zf`1.7,O,!~zb,<,?6`.7,6:O,`.5,5:?z,2:*/OH,3:/*N?!          ")
  %-----------------------------------------------------------------------
  putm
  **(defaultscale:=0.8;
      label.lrt("EN: "&EN,(0,h));
      label.lrt("fm: "&fm,(0,h-5mm));
      label.lrt("mw: "&mw,(0,h-9mm));
      label.lrt("MW: "&MW,(0,h-13mm));)
endfigm
\end{mplibcode}
%%%%%%%%%%%%%%%%%%%%%%%%%%%%%%%%%%%%%%%%%%%%%%%%%%%%%%%%%%%%%%%%%%%%%%%%%%%%%
\paragraph{(Paclitaxel)}
\noindent%
\begin{verbatim}
beginfigm
  % "EN:Paclitaxel","MW:853.918",
  readm(
  %------------------------------------------------------------------
  "?6,5=dl,3:,`1,36,45,45,45,45,`,:5,-4=?6,-4=?4,-1=wb,-3=wf,-1:O,  ",
  " 4:??,6:?,3^-60'15:*/OH,8:/*H^-60,9:?w^60,10:?O,                 ",
  " 1:,!,O,!,?O,!,*/OH,!,/Ph,60~wf,NH,-60,?O,60,Ph,                 ",
  " 7:,!z,O,-45,?O,60,Ph,11^rl'*12^-15^lr:*/O!'?O!                  ")
  %------------------------------------------------------------------
  #(140mm,30mm)
  @(0,0.5) putm
  sw_numbering:=Atom+All;
  @(0.6,0.5) putm
  sw_numbering:=Bond+All;
  @(1,0.5) putm
endfigm
\end{verbatim}
%----------------------------------------------------------------------------
\begin{mplibcode}
beginfigm
  % "EN:Paclitaxel","MW:853.918",
  readm(
  %-------------------------------------------------------------------
  "?6,5=dl,3:,`1,36,45,45,45,45,`,:5,-4=?6,-4=?4,-1=wb,-3=wf,-1:O,  ",
  " 4:??,6:?,3^-60'15:*/OH,8:/*H^-60,9:?w^60,10:?O,                 ",
  " 1:,!,O,!,?O,!,*/OH,!,/Ph,60~wf,NH,-60,?O,60,Ph,                 ",
  " 7:,!z,O,-45,?O,60,Ph,11^rl'*12^-15^lr:*/O!'?O!                  ")
  %-------------------------------------------------------------------
  #(160mm,40mm)
  @(0,  0.5)                         putm
  sw_numbering:=Atom+All; @(0.5,0.5) putm
  sw_numbering:=Bond+All; @(1,  0.5) putm
endfigm
\end{mplibcode}
%%%%%%%%%%%%%%%%%%%%%%%%%%%%%%%%%%%%%%%%%%%%%%%%%%%%%%%%%%%%%%%%%%%%%%%%%
\noindent%
\newpage
\subsection{loadm example}
\index{sw\_comment}%
\index{loadm()}%
\paragraph{(Example)}
\begin{verbatim}
loadm("CAT=biological","MW>=285","MW<=288","a:EN");
\end{verbatim}
%---------------------------------------------------------------
\paragraph{(output)}
\begin{verbatim}
* jobname=mcf_exa_soc
* numbersystem=double
* output report file
* file name=mcf_exa_soc-report.txt)
* mcf_template 2023.05.07
* Input  : main_lib.mcf [525]
* Output : ucount [4]
* Filter(1): CAT =biological
* Filter(2): MW >= 285
* Filter(3): MW <= 288
* Sort key : EN (ascending)
[1]:Luteolin
[2]:Lycorine
[3]:Morphine
[4]:Piperine )

row[1][1]="CAT:biological;EN:Luteolin;MW:286.24;LV:-"
row[1][2]=":"
row[1][3]="^30,Ph,3=?6,9=dl,10:O,7:?O,9:,!,Ph,2'6'14'15:/OH"
row[1][4]=";"
row[2][1]="CAT:biological;EN:Lycorine;MW:287.315;LV:1"
row[2][2]=":"
row[2][3]="^30,Ph,-4'-2=?6,6'9--12=?5,13=dl,8:N,15'17:O,"
row[2][4]="*9^180'10^60:*/H,13'*14:*/OH"
row[2][5]=";"
row[3][1]="CAT:biological;EN:Morphine;MW:285.343;LV:1"
row[3][2]=":"
row[3][3]="^30,Ph,2'-4=?6,1---12=?5,-1:O,-1=zb,"
row[3][4]="7:,60~wf`0.75,70~si_`1.3,45,N!,:9~wb,15=dl,6:/OH,8^180:*/H,12:/*OH"
row[3][5]=";"
row[4][1]="CAT:biological;EN:Piperine;MW:285.343;LV:1"
row[4][2]=":"
row[4][3]="^30,Ph,-1=?5,-1'-3:O,4:,!,!d,!,!d,!,?O,!,?6,-6:N"
row[4][4]=";"
%------------------------------------------------------------------------------
\end{verbatim}
\paragraph{(Tag)}
\begin{verbatim}
CAT : category                          USE : the use
EN  : english name                      JN  : japanese name
FM  : formula from data                 MW  : molecular weight from data
MI  : monoisotopic mass from data       CAS : CAS number
LV  : level                             LN  : long name
\end{verbatim}
\newpage
\noindent%
%----------------------------------------------------------------------------
\newpage
\subsection{getm example}
\paragraph{(Chlorophyll a)}
\noindent%
\begin{verbatim}
beginfigm
  getm("Chlorophyll a")
  sw_output:=Fig+Mcode;
  #(80mm,30mm)
  putm
  VerbatimTeX("\gdef\EN{"&EN&"}\gdef\MW{"&MW&"}");
  VerbatimTeX("\gdef\mw{"&mw&"}\gdef\fm{"&fm&"}");
endfigm
\end{mplibcode}
\verbatiminput{temp-mc.aux}            %%%% input temp-mc.aux %%%%
{\tt ** EN:\EN \quad mw:\MW \quad MW:\mw \quad fm:\fm}%
\end{verbatim}
%------------------------------------------------------------------------------------
\begin{mplibcode}
beginfigm
  sw_output:=Fig+Mcode;
  #(80mm,30mm)
  getm("Chlorophyll a") putm
  VerbatimTeX("\gdef\EN{"&EN&"}\gdef\MW{"&MW&"}");
  VerbatimTeX("\gdef\mw{"&mw&"}\gdef\fm{"&fm&"}");
endfigm
\end{mplibcode}
\verbatiminput{temp-mc.aux}            %%%% input temp-mc.aux %%%%
{\tt ** EN:\EN \quad mw:\MW \quad MW:\mw \quad fm:\fm}%
%%%%%%%%%%%%%%%%%%%%%%%%%%%%%%%%%%%%%%%%%%%%%%%%%%%%%%%%%%%%%%%%%%%%%%%%%%%%%%%%%%%%%
\paragraph{(Dinophysistoxin-1)}
\noindent%
\begin{verbatim}
beginfigm
  sw_output:=Fig+Mcode;               %%%% output temp-mc.aux   %%%%
  getm("Okadaic acid")                %%%% select Okadaic acid  %%%%
  readm(",40:?w,65=red")              %%%% add methyl group     %%%%
  EN:="Dinophysistoxin-1"; #(90mm,20mm) MW:="819";
  putm
  VerbatimTeX("\gdef\EN{"&EN&"}\gdef\MW{"&MW&"}");
  VerbatimTeX("\gdef\mw{"&mw&"}\gdef\fm{"&fm&"}");
endfigm
\end{mplibcode}
\verbatiminput{temp-mc.aux}            %%%% input temp-mc.aux %%%%
{\tt ** EN:\EN \quad mw:\MW \quad MW:\mw \quad fm:\fm}%
\end{verbatim}
%----------------------------------------------------------------------------
\begin{mplibcode}
beginfigm
  sw_output:=Fig+Mcode;          %%%% output temp-mc.aux   %%%%
  getm("Okadaic acid")           %%%% select Okadaic acid  %%%%
  readm(",40:?w,65=red")         %%%% add methyl group     %%%%
  MW:="819"; EN:="Dinophysistoxin-1"; #(90mm,20mm);
  putm
  VerbatimTeX("\gdef\EN{"&EN&"}\gdef\MW{"&MW&"}");
  VerbatimTeX("\gdef\mw{"&mw&"}\gdef\fm{"&fm&"}");
endfigm
\end{mplibcode}
\verbatiminput{temp-mc.aux}            %%%% input temp-mc.aux %%%%
{\tt ** EN:\EN \quad mw:\MW \quad MW:\mw \quad fm:\fm}%
%%%%%%%%%%%%%%%%%%%%%%%%%%%%%%%%%%%%%%%%%%%%%%%%%%%%%%%%%%%%%%%%%%%%%%%%%%%%%
\newpage
%----------------------------------------------------------------------------
\paragraph{(Maitotoxin)}
\noindent%
%--------------------------------------------------------------------------------
\begin{verbatim}
%--------------------------------------------------------------------------------
\begin{mplibcode}
  beginfigm
    getm("Maitotoxin") #(150mm,80mm) #@(3mm,3mm)
    sw_output:=Fig+Mcode;         %%%% output temp-mc.aux %%%%
    sw_frame:=Outside;
    putm
    VerbatimTeX("\gdef\EN{"&EN&"}\gdef\MW{"&MW&"}");
    VerbatimTeX("\gdef\mw{"&mw&"}\gdef\fm{"&fm&"}");
  endfigm
\end{mplibcode}
\verbatiminput{temp-mc.aux}            %%%% input temp-mc.aux %%%%
{\tt ** EN:\EN \quad mw:\MW \quad MW:\mw \quad fm:\fm}%
%--------------------------------------------------------------------------------
\end{verbatim}
%--------------------------------------------------------------------------------
\begin{mplibcode}
  beginfigm
    getm("Maitotoxin") #(150mm,80mm) #@(3mm,3mm) 
    sw_output:=Fig+Mcode;
    sw_frame:=Outside; %% mc_length:=40;
    putm
    VerbatimTeX("\gdef\EN{"&EN&"}\gdef\MW{"&MW&"}");
    VerbatimTeX("\gdef\mw{"&mw&"}\gdef\fm{"&fm&"}");
  endfigm
\end{mplibcode}
\verbatiminput{temp-mc.aux}
{\tt ** EN:\EN \quad mw:\MW \quad MW:\mw \quad fm:\fm}%
%--------------------------------------------------------------------------------
%%%%%%%%%%%%%%%%%%%%%%%%%%%%%%%%%%%%%%%%%%%%%%%%%%%%%%%%%%%%%%%%%%%%%%%%%%%%%
%------------------------------------------------------------------------------
\newpage
\section{Example to use mcf2graph}
\subsection{MetaPost souce file}
\index{mcf2graph.mp}%
\index{sw\_output}%
\index{\textbackslash()}%
\index{plus}%
\index{getm()}%
\index{putm}%
\index{ucount}%
\begin{verbatim}
input mcf2graph;                                % input main macro
#(60mm,40mm);                                   %  (figure width,figure height)         %
outputformat:="png"; hppp:=vppp:=0.1;           % PNG output
outputtemplate:="c%3c-%{EN_}.png";              %
%------------------------------------------------------------------------------
beginfigm
  % EN:Ampicillin  MW:349.405
  drawm("^45,?4,-3=?5,2:N,7:S,"&                        % immediately write
    "3^45:/*H,1:?O^15,5:/*COOH^-18,6:??,"&
    "4:,!w^15,NH,!,?O,!,/*NH2,!,Ph")
endfigm
%------------------------------------------------------------------------------
beginfigm
  % EN:Cholesterol  MW:386.65
  readm("^30,?6,-4'-2=?6,-4=?5,7=dl,",                        % mc1
        "10:/*H^180,11:/*H^-60,17:/*H^-54,4'12:?w^60, ",      % mc2
        "-1:,18,?z,-60,!3,?!")                                % mc3
  putm                                                        % put figure
endfigm
%------------------------------------------------------------------------------
loadm()                                                 % load all unit
beginfigm getm("Adenine") putm endfigm                  % get EN=Adenine
beginfigm getm(4) putm endfigm                          % select No.4
for i=1 upto ucount:                                    % figure count
  beginfigm getm(i) putm endfigm                        % select No.i
endfor
bye
\end{verbatim}
%------------------------------------------------------------------------
\noindent%
\subsection{Molecular library file}
\begin{verbatim}
% first character of line
% "%":comment out / ":": start MCF / ";": stop MCF
% "=":start parameter setting /  "*": start ext(...) / "+": start add(...)
% CAT : Category,EN : Name,MW : Molecular weight LV : Level
%------------------------------------------------------------------------
CAT:-;EN:-;MW:0;LV:-
=
 sw_frame:=Atom;
:
^30,?6,3=?5,1'3'5'9=dl,2'6'9:N,5:/NH2,7:NH
*
 defaultscale:=.5; label.bot(decimal(fig_num)&":"&EN,(.5w,0));
+
 defaultscale:=.3; label.bot("A2",A2) withcolor red;
;
%==============================================================================
CAT:biological;EN:Adenine;MW:135.13;LV:1
:
^30,?6,3=?5,1'3'5'9=dl,2'6'9:N,5:/NH2,7:NH
;
%------------------------------------------------------------------------------
CAT:biological;EN:Guanine;MW:151.13;LV:1
:
^30,?6,3=?5,1'3'9=dl,2'9:N,6'7:NH,5:?O,1:/NH2
;
%------------------------------------------------------------------------------
\end{verbatim}
%------------------------------------------------------------------------
\noindent%
\newpage
\subsection{MCF aux file output}
\paragraph{(Option parameter setting)}
\index{Mcode}%
\begin{verbatim}
sw_output:=Mcode;         %% output 'temp-mc.aux'
\end{verbatim}
\paragraph{(Command line)}
\begin{verbatim}
  >mpost -s ahlength=3 FILENAME  (sw_output=Fig Expand mode)
\end{verbatim}
\paragraph{(Output mcf file)}
\begin{verbatim}
sw_output=Mcode           %% file name = 'temp-mc.aux'

(result)
^30,?6,3=?5,1'3'5'9=dl,2'6'9:N,5:/NH2,7:NH

\end{verbatim}
%------------------------------------------------------------------------
\paragraph{(LuaLaTeX example)}
%-----------------------------------------------------------------------
\begin{verbatim}
%-----------------------------------------------------------------------
%%  "EN:Vancomycin
\begin{mplibcode}
  beginfigm
    sw_output:=Mcode;     %%%% output temp-mc.aux %%%%
  endfigm
\end{mplibcode}
%-----------------------------------------------------------------------
\verbatiminput{temp-mc.aux}
%-----------------------------------------------------------------------
\end{verbatim}
%-----------------------------------------------------------------------
\begin{verbatim}
(result)
file name = 'temp-mc.aux'

^30,?6,4:,?6,-4:,!,!3,^-12,?5,-3:@,^-12,?6,-3=?6,-3:,!w,!3,
 ?6,-4:,?6,6:,!,!,?z^-40,*/OH^20,!,?O,!1,OH,
 3=wb,11=dl,15=dr,17'19=wf,38=wb,5'7'16'24'25'33'42:O,
 32:*/H^60,10:?,12'31:?w,27:?d,37:?z,28:/OH,3'29:/*OH
\end{verbatim}
%------------------------------------------------------------------------
\newpage
\noindent%
\subsection{Report output}
\paragraph{(Option parameter setting)}
\index{sw\_output}%
\index{Report}%
\begin{verbatim}
  sw_output:=Report;        %% file name = 'jobname-report.aux'
\end{verbatim}
\paragraph{(Command line)}
\begin{verbatim}
  >mpost -s ahlength=7 FILENAME
\end{verbatim}
\paragraph{(Output)}
\begin{verbatim}
===========================================================================
 No[148],Name<Phenol>,Category<synthetic>,File<main_lib.mcf>
---------------------------------------------------------------------------
 Row[1],Length[12],Block[3],BackboneA[6],BackboneB[6],Group[1]
---------------------------------------------------------------------------
^30,Ph,3:/OH
---------------------------------------------------------------------------
[1  ]  ^30
[2  ]  Ph
[3  ]  3:/OH
---------------------------------------------------------------------------
---------------------------------------------------------------------------
 Width[35.80607],Height[24.55503], Shift x[-1.77635],Shift y[-7.54719]
 Bond length[11.33855],Atom size[4.881881]
---------------------------------------------------------------------------
 Atom[7],Bond[7],Ring[1],Hide H[5]
< NO. ><atom(s) >(  x axis   ,   y axis   )<bond><hideH><chg>
 A1     C        (         0 ,          0 )    3     1
 A2     C        (         1 ,         -1 )    3     1
 A3     C        (         2 ,          0 )    4        
 A4     C        (         2 ,          1 )    3     1
 A5     C        (         1 ,          1 )    3     1
 A6     C        (         0 ,          1 )    3     1
 A7     OH       (         3 ,          0 )    1        
---------------------------------------------------------------------------
< NO. ><  bond   (sdt)><angle +(  +-  )><length (   pt   )>
 B1     1 ->   2 (  2)     330 (   -30)       1 (      11)
 B2     2 ->   3 (  1)      30 (    30)       1 (      11)
 B3     3 ->   4 (  2)      90 (    90)       1 (      11)
 B4     4 ->   5 (  1)     150 (   150)       1 (      11)
 B5     5 ->   6 (  2)     210 (  -150)       1 (      11)
 B6     6 ->   1 (  1)     270 (   -90)       1 (      11)
 B7     3 ->   7 (  1)     330 (   -30)0.660000 (       7)
---------------------------------------------------------------------------
<atom>( atom wt )[ mi wt   ]  < cnt > < sum wt   >[ sum mi wt  ]
 C    (  12.0107)[       12] *        6    72.0642[          72]
 H    (1.0079400)[1.0078250] *        6    6.04764[6.0469501933]
 O    (  15.9994)[15.994914] *        1    15.9994[15.994914619]
 Molecular Weight [Mono Isotopic] =       94.11123[  94.0418648]
---------------------------------------------------------------------------
 Weight Calc: 94.11123 - Input: 94.11 = 0.0012399
 Fomula Calc: C6H6O
===========================================================================
\end{verbatim}
\newpage
%------------------------------------------------------------------------
\noindent%
\subsection{MOL file output}
\paragraph{(Option parameter setting)}
\index{sw\_output}%
\index{MOL2000}%
\index{MOL3000}%
\begin{verbatim}
  sw_output:=MOL2000;   % MOL(V2000)
  sw_output:=MOL3000;   % MOL(V3000)
\end{verbatim}
\paragraph{(Command line)}
\begin{verbatim}
  >mpost -s ahlength=5  FILENAME     % MOL(V2000)
  >mpost -s ahlength=6  FILENAME     % MOL(V3000)
\end{verbatim}
\paragraph{(Output)}
\begin{verbatim}
%%%%%%%%%%%%%%%%%%%%%%%%%%%%%%%%%%%%%%%%%%%%%%%%%%%%%%%%%%%%%%%%%%%%%
  -MCFtoMOL- EN:Caffeine         

 14 15  0  0  0  0  0  0  0  0999 V2000
         0         0         0 C   0  0  0  0
   0.86603      -0.5         0 N   0  0  0  0
   1.73206         0         0 C   0  0  0  0
   1.73206         1         0 C   0  0  0  0
   0.86603       1.5         0 C   0  0  0  0
         0         1         0 N   0  0  0  0
    2.6831  -0.30902         0 N   0  0  0  0
   3.27089       0.5         0 C   0  0  0  0
    2.6831   1.30902         0 N   0  0  0  0
   0.86603  -1.36383         0 C   0  0  0  0
  -0.76894   1.44394         0 C   0  0  0  0
  -0.76894  -0.44394         0 O   0  0  0  0
   0.86603   2.36383         0 O   0  0  0  0
   2.95299    2.1396         0 C   0  0  0  0
  1  2  1  0     0  0
  2  3  1  0     0  0
  3  4  2  0     0  0
  4  5  1  0     0  0
  5  6  1  0     0  0
  6  1  1  0     0  0
  3  7  1  0     0  0
  7  8  2  0     0  0
  8  9  1  0     0  0
  9  4  1  0     0  0
  2 10  1  0     0  0
  6 11  1  0     0  0
  1 12  2  0     0  0
  5 13  2  0     0  0
  9 14  1  0     0  0
M  END
%%%%%%%%%%%%%%%%%%%%%%%%%%%%%%%%%%%%%%%%%%%%%%%%%%%%%%%%%%%%%%%%%%%%%
\end{verbatim}
%----------------------------------------------------------------------------
\newpage
\subsection{LuaTeX file example}
%############################################################################
\begin{verbatim}
** file mcf2graph.mp must be in the same directory to typeset

\documentclass{article}
\usepackage{luamplib}%
\mplibcodeinherit{enable}%
\mplibverbatim{enable}%
\mplibnumbersystem{double}%
\begin{document}
\noindent%
%--------------------------------------------------------------------
\begin{mplibcode}
  input mcf2graph;
  max_blength:=4.5mm;
  defaultsize:=8bp;
  defaultscale:=1;
  %--------------------------------------------------------
  EN:="Limonin";
  MW="470.51";
  beginfigm
    #(50mm,50mm)
    readm(
    %---------------------------------------
    "^30,?6,-3'-4=?6,                     ",
    " -5=?3,-2=wf,-1=wb,6=?5,-4=?6,-5=wf, ",
    " 13'15'17'20:O,3'12'21:?O,           ",
    " 4~wf^60'8~zf^60'18^35'18^-35:?,     ",
    " 1^60'5^180'16^60:/*H,               ",
    " :14,!z,<,?5,1'4=dl,3:O              ")
    %---------------------------------------
    putm
  endfigm
\end{mplibcode}\\
%--------------------------------------------------------------------
\begin{mplibcode}
  EN:="beta-carotene";
  MW:="536.87";
  beginfigm
    #(80mm,50mm)
    readm(
    %-----------------------------------------
    "^30,?6,3=dl,3'5^35'5^-35:?,           ",
    " :4,!,<,!18,1'3'5'7'9'11'13'15'17=dr, ",
    " 3'7'12'16:?,                         ",
    " <,?6,6=dl,6'2^35'2^-35:?             ")
    %-----------------------------------------
    putm
  endfigm
\end{mplibcode}\\
%--------------------------------------------------------------------
\end{document}
\end{verbatim}
%############################################################################
\newpage
\subsection{Batch file mcf\_to\_graph.cmd}
\index{mcf\_to\_graph.cmd}%
%############################################################################
Batch file for compile mcf or mp file
\begin{verbatim}
** drag and drop mcf or mp file on this batch : >mcf_to_graph filenam
** file mcf2graph.mp , main_lib.mcf must be in the same directory

:::::::::::::::::::::::::::::::::::::::::::::::::::::::::::::::::::::::::::::::
::  Batch file for compile MCF  2025.09.15
::  <drag and drop library or metapost file on this batch>
:::::::::::::::::::::::::::::::::::::::::::::::::::::::::::::::::::::::::::::::
@echo off
if %~x1 == .mp (
if %~n1 == mcf2graph goto end
mkdir %~n1-mp
mpost -output-directory=./%~n1-mp %~n1
goto end )
if not %~x1 == .mcf goto end
echo **** Select output format / Library file [%~n1] ****
echo [1:svg 2:png-600dpi 3:png-1200dpi 4:MOL-v2000 5:MOL-v3000 6:report 0:cancel]
choice /c 1234560
if %errorlevel% == 7 goto end
echo input mcf2graph; libm:="%~n1.mcf"; loadm() allm bye> temp_soc.mp
if %errorlevel% == 1 (mkdir %~n1-svg
mpost -output-directory=./%~n1-svg temp_soc)
if %errorlevel% == 2 (mkdir %~n1-pn0600
mpost -output-directory=./%~n1-pn0600 -s ahangle=1 temp_soc)
if %errorlevel% == 3 (mkdir %~n1-pn1200
mpost -output-directory=./%~n1-pn1200 -s ahangle=2 temp_soc)
if %errorlevel% == 4 (mkdir %~n1-molv2k
mpost -output-directory=./%~n1-molv2k -s ahlength=5 temp_soc)
if %errorlevel% == 5 (mkdir %~n1-molv3k
mpost -output-directory=./%~n1-molv3k -s ahlength=6 temp_soc)
if %errorlevel% == 6 (mkdir %~n1-report
mpost -output-directory=./%~n1-report -numbersystem=double -s ahlength=7 temp_soc)
del temp_soc.mp
:end
\end{verbatim}
%############################################################################
%------------------------------------------------------------------------
\texttt{\printindex}
%------------------------------------------------------------------------
\end{document}
----------------------
\end{document}
